% Chinese adaptation of this SGA 4 source's own AMS/SMF layout.
\usepackage{xeCJK}
\setCJKmainfont[BoldFont={Noto Serif CJK SC Bold},ItalicFont={Noto Serif CJK SC},SmallCapsFont={Noto Serif CJK SC}]{Noto Serif CJK SC}
\setCJKsansfont{Noto Sans CJK SC}
\setCJKmonofont{Noto Sans Mono CJK SC}
\def\chaptername{第}
\def\contentsname{目录}
\def\bibname{参考文献}
\def\appendixname{附录}
\def\proofname{证明}
\def\theoname{定理}\def\propname{命题}\def\conjname{猜想}
\def\coroname{推论}\def\lemmname{引理}\def\definame{定义}
\def\remaname{注}\def\exemname{例}
\def\cf{\ifmmode\text{参见}\else 参见\fi\kern.3em}\def\Cf{参见\kern.3em}
\def\ie{\ifmmode\text{即}\else 即\fi\kern.3em}\def\iev{即\xspace}
\def\resp{\ifmmode\text{分别}\else 分别\fi}\def\etcn{等等\xspace}\def\etco{等等\xspace}
\def\quaet{\quad\text{且}\quad}
\def\loccit{前引文献\xspace}
% Heading label follows original Roman exposé numbering.
\def\thechapter{\Roman{chapter}}
% SGA4 uses numbered paragraphs below subsubsections. The inherited AMS class
% stops at level3, suppressing those printed numbers despite explicit labels.
\setcounter{secnumdepth}{5}
\hypersetup{pdftitle={SGA 4：拓扑斯理论与概形的平展上同调，第 X--XIX 讲（中文首译稿）},pdfauthor={M. Artin; A. Grothendieck; J.-L. Verdier},pdfsubject={SGA 4 中文原生 LaTeX，覆盖状况见 PROGRESS.json},pdfkeywords={SGA 4, 拓扑斯, 平展上同调, 中文翻译, X--XIX}}

\emergencystretch=2em
\usepackage{etoolbox}
\makeatletter
\patchcmd{\@makechapterhead}{\enspace\thechapter}{\enspace\thechapter\enspace 讲}{}{}
\makeatother
\AtBeginDocument{\def\chaptername{第}\def\contentsname{目录}\def\bibname{参考文献}\def\proofname{证明}}
% Generated source-backed external reference registry. No synthetic content/pages.
\expandafter\def\csname SGAsource@5.3'\endcsname{3'}
\expandafter\def\csname SGAsource@5.4'\endcsname{4'}
\expandafter\def\csname SGAsource@I.4.3.1\endcsname{4.3.1}
\expandafter\def\csname SGAsource@I.8.12.1.1\endcsname{8.12.1.1}
\expandafter\def\csname SGAsource@I.8.9.12\endcsname{8.9.12}
\expandafter\def\csname SGAsource@I.9.12.1\endcsname{9.12.1}
\expandafter\def\csname SGAsource@I.condef1.3\endcsname{1.3}
\expandafter\def\csname SGAsource@I.cor2.10\endcsname{2.10}
\expandafter\def\csname SGAsource@I.cor2.3.2\endcsname{2.3.2}
\expandafter\def\csname SGAsource@I.cor2.9\endcsname{2.9}
\expandafter\def\csname SGAsource@I.cor3.2\endcsname{3.2}
\expandafter\def\csname SGAsource@I.cor3.3\endcsname{3.3}
\expandafter\def\csname SGAsource@I.cor3.5\endcsname{3.5}
\expandafter\def\csname SGAsource@I.cor5.5.1\endcsname{5.5.1}
\expandafter\def\csname SGAsource@I.cor6.3\endcsname{6.3}
\expandafter\def\csname SGAsource@I.cor7.11\endcsname{7.11}
\expandafter\def\csname SGAsource@I.cor7.12\endcsname{7.12}
\expandafter\def\csname SGAsource@I.cor7.3\endcsname{7.3}
\expandafter\def\csname SGAsource@I.cor7.5\endcsname{7.5}
\expandafter\def\csname SGAsource@I.cor7.7.1\endcsname{7.7.1}
\expandafter\def\csname SGAsource@I.cor8.1.7\endcsname{8.1.7}
\expandafter\def\csname SGAsource@I.cor8.11.3\endcsname{8.11.3}
\expandafter\def\csname SGAsource@I.cor8.12.7\endcsname{8.12.7}
\expandafter\def\csname SGAsource@I.cor8.12.8\endcsname{8.12.8}
\expandafter\def\csname SGAsource@I.cor8.13.3\endcsname{8.13.3}
\expandafter\def\csname SGAsource@I.cor8.13.4\endcsname{8.13.4}
\expandafter\def\csname SGAsource@I.cor8.5.2\endcsname{8.5.2}
\expandafter\def\csname SGAsource@I.cor8.6.5\endcsname{8.6.5}
\expandafter\def\csname SGAsource@I.cor8.7.6\endcsname{8.7.6}
\expandafter\def\csname SGAsource@I.cor8.7.7\endcsname{8.7.7}
\expandafter\def\csname SGAsource@I.cor8.9.2\endcsname{8.9.2}
\expandafter\def\csname SGAsource@I.cor8.9.3\endcsname{8.9.3}
\expandafter\def\csname SGAsource@I.cor8.9.8\endcsname{8.9.8}
\expandafter\def\csname SGAsource@I.cor9.15\endcsname{9.15}
\expandafter\def\csname SGAsource@I.cor9.17\endcsname{9.17}
\expandafter\def\csname SGAsource@I.cor9.18\endcsname{9.18}
\expandafter\def\csname SGAsource@I.cor9.19\endcsname{9.19}
\expandafter\def\csname SGAsource@I.cor9.20\endcsname{9.20}
\expandafter\def\csname SGAsource@I.cor9.21.10.1\endcsname{9.21.10.1}
\expandafter\def\csname SGAsource@I.cor9.23\endcsname{9.23}
\expandafter\def\csname SGAsource@I.cor9.24\endcsname{9.24}
\expandafter\def\csname SGAsource@I.cor9.25\endcsname{9.25}
\expandafter\def\csname SGAsource@I.cor9.26\endcsname{9.26}
\expandafter\def\csname SGAsource@I.cor9.7\endcsname{9.7}
\expandafter\def\csname SGAsource@I.cor9.8\endcsname{9.8}
\expandafter\def\csname SGAsource@I.cor9.9\endcsname{9.9}
\expandafter\def\csname SGAsource@I.def1.1\endcsname{1.1}
\expandafter\def\csname SGAsource@I.def1.2\endcsname{1.2}
\expandafter\def\csname SGAsource@I.def2.1\endcsname{2.1}
\expandafter\def\csname SGAsource@I.def2.2\endcsname{2.2}
\expandafter\def\csname SGAsource@I.def2.4.1\endcsname{2.4.1}
\expandafter\def\csname SGAsource@I.def2.5\endcsname{2.5}
\expandafter\def\csname SGAsource@I.def2.7\endcsname{2.7}
\expandafter\def\csname SGAsource@I.def4.1\endcsname{4.1}
\expandafter\def\csname SGAsource@I.def6.1\endcsname{6.1}
\expandafter\def\csname SGAsource@I.def7.1\endcsname{7.1}
\expandafter\def\csname SGAsource@I.def8.1.1\endcsname{8.1.1}
\expandafter\def\csname SGAsource@I.def8.1.8\endcsname{8.1.8}
\expandafter\def\csname SGAsource@I.def9.1\endcsname{9.1}
\expandafter\def\csname SGAsource@I.def9.12\endcsname{9.12}
\expandafter\def\csname SGAsource@I.def9.2\endcsname{9.2}
\expandafter\def\csname SGAsource@I.def9.3\endcsname{9.3}
\expandafter\def\csname SGAsource@I.def9.4\endcsname{9.4}
\expandafter\def\csname SGAsource@I.eq7.6.1\endcsname{1}
\expandafter\def\csname SGAsource@I.exem2.7.1\endcsname{2.7.1}
\expandafter\def\csname SGAsource@I.exem7.13\endcsname{7.13}
\expandafter\def\csname SGAsource@I.exem8.11.6\endcsname{8.11.6}
\expandafter\def\csname SGAsource@I.exem9.12.2\endcsname{9.12.2}
\expandafter\def\csname SGAsource@I.exer8.12.9\endcsname{8.12.9}
\expandafter\def\csname SGAsource@I.exer8.2.8\endcsname{8.2.8}
\expandafter\def\csname SGAsource@I.exer8.5.5\endcsname{8.5.5}
\expandafter\def\csname SGAsource@I.exer8.7.8\endcsname{8.7.8}
\expandafter\def\csname SGAsource@I.exer8.7.9\endcsname{8.7.9}
\expandafter\def\csname SGAsource@I.exer8.8.6\endcsname{8.8.6}
\expandafter\def\csname SGAsource@I.exer8.8.7\endcsname{8.8.7}
\expandafter\def\csname SGAsource@I.exer8.8.8\endcsname{8.8.8}
\expandafter\def\csname SGAsource@I.exer8.9.10\endcsname{8.9.10}
\expandafter\def\csname SGAsource@I.exer8.9.11\endcsname{8.9.11}
\expandafter\def\csname SGAsource@I.exer8.9.9\endcsname{8.9.9}
\expandafter\def\csname SGAsource@I.exer9.20.1\endcsname{9.20.1}
\expandafter\def\csname SGAsource@I.lem2.8.1\endcsname{2.8.1}
\expandafter\def\csname SGAsource@I.lem6.5.2\endcsname{6.5.2}
\expandafter\def\csname SGAsource@I.lem7.6.2\endcsname{7.6.2}
\expandafter\def\csname SGAsource@I.lem8.12.4\endcsname{8.12.4}
\expandafter\def\csname SGAsource@I.lem8.12.8.1\endcsname{8.12.8.1}
\expandafter\def\csname SGAsource@I.lem8.8.3.4\endcsname{8.8.3.4}
\expandafter\def\csname SGAsource@I.lem8.8.5.1\endcsname{8.8.5.1}
\expandafter\def\csname SGAsource@I.lem9.11.1\endcsname{9.11.1}
\expandafter\def\csname SGAsource@I.lem9.11.2\endcsname{9.11.2}
\expandafter\def\csname SGAsource@I.lem9.13.1\endcsname{9.13.1}
\expandafter\def\csname SGAsource@I.lem9.16.2\endcsname{9.16.2}
\expandafter\def\csname SGAsource@I.lem9.21.10\endcsname{9.21.10}
\expandafter\def\csname SGAsource@I.lem9.21.13\endcsname{9.21.13}
\expandafter\def\csname SGAsource@I.lem9.21.16\endcsname{9.21.16}
\expandafter\def\csname SGAsource@I.lem9.21.17\endcsname{9.21.17}
\expandafter\def\csname SGAsource@I.lem9.21.18\endcsname{9.21.18}
\expandafter\def\csname SGAsource@I.lem9.21.19\endcsname{9.21.19}
\expandafter\def\csname SGAsource@I.lem9.21.9\endcsname{9.21.9}
\expandafter\def\csname SGAsource@I.not5.9\endcsname{5.9}
\expandafter\def\csname SGAsource@I.prop1.4\endcsname{1.4}
\expandafter\def\csname SGAsource@I.prop2.11\endcsname{2.11}
\expandafter\def\csname SGAsource@I.prop2.12\endcsname{2.12}
\expandafter\def\csname SGAsource@I.prop2.3\endcsname{2.3}
\expandafter\def\csname SGAsource@I.prop2.3.3\endcsname{2.3.3}
\expandafter\def\csname SGAsource@I.prop2.6\endcsname{2.6}
\expandafter\def\csname SGAsource@I.prop2.8\endcsname{2.8}
\expandafter\def\csname SGAsource@I.prop3.1\endcsname{3.1}
\expandafter\def\csname SGAsource@I.prop3.4\endcsname{3.4}
\expandafter\def\csname SGAsource@I.prop3.6\endcsname{3.6}
\expandafter\def\csname SGAsource@I.prop4.2\endcsname{4.2}
\expandafter\def\csname SGAsource@I.prop5.1\endcsname{5.1}
\expandafter\def\csname SGAsource@I.prop5.11\endcsname{5.11}
\expandafter\def\csname SGAsource@I.prop5.14\endcsname{5.14}
\expandafter\def\csname SGAsource@I.prop5.2\endcsname{5.2}
\expandafter\def\csname SGAsource@I.prop5.4\endcsname{5.4}
\expandafter\def\csname SGAsource@I.prop5.5\endcsname{5.5}
\expandafter\def\csname SGAsource@I.prop5.6\endcsname{5.6}
\expandafter\def\csname SGAsource@I.prop5.8.3\endcsname{5.8.3}
\expandafter\def\csname SGAsource@I.prop6.2\endcsname{6.2}
\expandafter\def\csname SGAsource@I.prop6.4\endcsname{6.4}
\expandafter\def\csname SGAsource@I.prop6.5\endcsname{6.5}
\expandafter\def\csname SGAsource@I.prop7.10\endcsname{7.10}
\expandafter\def\csname SGAsource@I.prop7.2\endcsname{7.2}
\expandafter\def\csname SGAsource@I.prop7.4\endcsname{7.4}
\expandafter\def\csname SGAsource@I.prop7.6\endcsname{7.6}
\expandafter\def\csname SGAsource@I.prop7.7\endcsname{7.7}
\expandafter\def\csname SGAsource@I.prop7.8\endcsname{7.8}
\expandafter\def\csname SGAsource@I.prop8.1.3\endcsname{8.1.3}
\expandafter\def\csname SGAsource@I.prop8.1.6\endcsname{8.1.6}
\expandafter\def\csname SGAsource@I.prop8.11.2\endcsname{8.11.2}
\expandafter\def\csname SGAsource@I.prop8.11.4\endcsname{8.11.4}
\expandafter\def\csname SGAsource@I.prop8.12.2\endcsname{8.12.2}
\expandafter\def\csname SGAsource@I.prop8.12.5\endcsname{8.12.5}
\expandafter\def\csname SGAsource@I.prop8.13.2\endcsname{8.13.2}
\expandafter\def\csname SGAsource@I.prop8.3.1\endcsname{8.3.1}
\expandafter\def\csname SGAsource@I.prop8.5.1\endcsname{8.5.1}
\expandafter\def\csname SGAsource@I.prop8.6.3\endcsname{8.6.3}
\expandafter\def\csname SGAsource@I.prop8.6.4\endcsname{8.6.4}
\expandafter\def\csname SGAsource@I.prop8.7.3\endcsname{8.7.3}
\expandafter\def\csname SGAsource@I.prop8.7.5\endcsname{8.7.5}
\expandafter\def\csname SGAsource@I.prop8.8.2\endcsname{8.8.2}
\expandafter\def\csname SGAsource@I.prop8.8.3\endcsname{8.8.3}
\expandafter\def\csname SGAsource@I.prop8.8.5\endcsname{8.8.5}
\expandafter\def\csname SGAsource@I.prop8.9.1\endcsname{8.9.1}
\expandafter\def\csname SGAsource@I.prop8.9.5\endcsname{8.9.5}
\expandafter\def\csname SGAsource@I.prop8.9.7\endcsname{8.9.7}
\expandafter\def\csname SGAsource@I.prop9.11\endcsname{9.11}
\expandafter\def\csname SGAsource@I.prop9.13\endcsname{9.13}
\expandafter\def\csname SGAsource@I.prop9.14\endcsname{9.14}
\expandafter\def\csname SGAsource@I.prop9.16\endcsname{9.16}
\expandafter\def\csname SGAsource@I.prop9.5\endcsname{9.5}
\expandafter\def\csname SGAsource@I.prop9.6\endcsname{9.6}
\expandafter\def\csname SGAsource@I.rem1.1.2\endcsname{1.1.2}
\expandafter\def\csname SGAsource@I.rem4.2.1\endcsname{4.2.1}
\expandafter\def\csname SGAsource@I.rem5.5.2\endcsname{5.5.2}
\expandafter\def\csname SGAsource@I.rem5.7\endcsname{5.7}
\expandafter\def\csname SGAsource@I.rem8.12.6\endcsname{8.12.6}
\expandafter\def\csname SGAsource@I.rem8.3.4\endcsname{8.3.4}
\expandafter\def\csname SGAsource@I.rem8.7.5.2\endcsname{8.7.5.2}
\expandafter\def\csname SGAsource@I.rem8.8.4\endcsname{8.8.4}
\expandafter\def\csname SGAsource@I.rem8.9.4\endcsname{8.9.4}
\expandafter\def\csname SGAsource@I.rem8.9.6\endcsname{8.9.6}
\expandafter\def\csname SGAsource@I.rem9.10\endcsname{9.10}
\expandafter\def\csname SGAsource@I.rem9.11.3\endcsname{9.11.3}
\expandafter\def\csname SGAsource@I.rem9.13.2\endcsname{9.13.2}
\expandafter\def\csname SGAsource@I.sec0\endcsname{0}
\expandafter\def\csname SGAsource@I.sec1\endcsname{1}
\expandafter\def\csname SGAsource@I.sec1.0\endcsname{1.0}
\expandafter\def\csname SGAsource@I.sec1.1.1\endcsname{1.1.1}
\expandafter\def\csname SGAsource@I.sec1.3.1\endcsname{1.3.1}
\expandafter\def\csname SGAsource@I.sec1.3.2\endcsname{1.3.2}
\expandafter\def\csname SGAsource@I.sec1.3.3\endcsname{1.3.3}
\expandafter\def\csname SGAsource@I.sec1.4.1\endcsname{1.4.1}
\expandafter\def\csname SGAsource@I.sec1.4.2\endcsname{1.4.2}
\expandafter\def\csname SGAsource@I.sec10\endcsname{10}
\expandafter\def\csname SGAsource@I.sec10.1\endcsname{10.1}
\expandafter\def\csname SGAsource@I.sec10.10\endcsname{10.10}
\expandafter\def\csname SGAsource@I.sec10.11\endcsname{10.11}
\expandafter\def\csname SGAsource@I.sec10.2\endcsname{10.2}
\expandafter\def\csname SGAsource@I.sec10.3\endcsname{10.3}
\expandafter\def\csname SGAsource@I.sec10.4\endcsname{10.4}
\expandafter\def\csname SGAsource@I.sec10.5\endcsname{10.5}
\expandafter\def\csname SGAsource@I.sec10.6\endcsname{10.6}
\expandafter\def\csname SGAsource@I.sec10.7\endcsname{10.7}
\expandafter\def\csname SGAsource@I.sec10.8\endcsname{10.8}
\expandafter\def\csname SGAsource@I.sec10.9\endcsname{10.9}
\expandafter\def\csname SGAsource@I.sec2\endcsname{2}
\expandafter\def\csname SGAsource@I.sec2.1.1\endcsname{2.1.1}
\expandafter\def\csname SGAsource@I.sec2.2.1\endcsname{2.2.1}
\expandafter\def\csname SGAsource@I.sec2.3.1\endcsname{2.3.1}
\expandafter\def\csname SGAsource@I.sec2.4.2\endcsname{2.4.2}
\expandafter\def\csname SGAsource@I.sec2.5.0\endcsname{2.5.0}
\expandafter\def\csname SGAsource@I.sec3\endcsname{3}
\expandafter\def\csname SGAsource@I.sec3.4.0\endcsname{3.4.0}
\expandafter\def\csname SGAsource@I.sec4\endcsname{4}
\expandafter\def\csname SGAsource@I.sec4.3\endcsname{4.3}
\expandafter\def\csname SGAsource@I.sec4.3.2\endcsname{4.3.2}
\expandafter\def\csname SGAsource@I.sec4.3.3\endcsname{4.3.3}
\expandafter\def\csname SGAsource@I.sec5\endcsname{5}
\expandafter\def\csname SGAsource@I.sec5.0\endcsname{5.0}
\expandafter\def\csname SGAsource@I.sec5.10\endcsname{5.10}
\expandafter\def\csname SGAsource@I.sec5.12\endcsname{5.12}
\expandafter\def\csname SGAsource@I.sec5.13\endcsname{5.13}
\expandafter\def\csname SGAsource@I.sec5.3\endcsname{5.3}
\expandafter\def\csname SGAsource@I.sec5.8.0\endcsname{5.8.0}
\expandafter\def\csname SGAsource@I.sec5.8.2\endcsname{5.8.2}
\expandafter\def\csname SGAsource@I.sec6\endcsname{6}
\expandafter\def\csname SGAsource@I.sec6.1.1\endcsname{6.1.1}
\expandafter\def\csname SGAsource@I.sec6.5.1\endcsname{6.5.1}
\expandafter\def\csname SGAsource@I.sec7\endcsname{7}
\expandafter\def\csname SGAsource@I.sec7.1.1\endcsname{7.1.1}
\expandafter\def\csname SGAsource@I.sec7.1.2\endcsname{7.1.2}
\expandafter\def\csname SGAsource@I.sec7.1.3\endcsname{7.1.3}
\expandafter\def\csname SGAsource@I.sec7.1.3.1\endcsname{7.1.3.1}
\expandafter\def\csname SGAsource@I.sec7.5.1\endcsname{7.5.1}
\expandafter\def\csname SGAsource@I.sec7.5.3\endcsname{7.5.3}
\expandafter\def\csname SGAsource@I.sec7.9\endcsname{7.9}
\expandafter\def\csname SGAsource@I.sec8\endcsname{8}
\expandafter\def\csname SGAsource@I.sec8.1\endcsname{8.1}
\expandafter\def\csname SGAsource@I.sec8.1.2\endcsname{8.1.2}
\expandafter\def\csname SGAsource@I.sec8.1.2.2\endcsname{8.1.2.2}
\expandafter\def\csname SGAsource@I.sec8.1.4\endcsname{8.1.4}
\expandafter\def\csname SGAsource@I.sec8.1.5\endcsname{8.1.5}
\expandafter\def\csname SGAsource@I.sec8.10\endcsname{8.10}
\expandafter\def\csname SGAsource@I.sec8.10.10\endcsname{8.10.10}
\expandafter\def\csname SGAsource@I.sec8.10.14\endcsname{8.10.14}
\expandafter\def\csname SGAsource@I.sec8.10.15\endcsname{8.10.15}
\expandafter\def\csname SGAsource@I.sec8.11\endcsname{8.11}
\expandafter\def\csname SGAsource@I.sec8.11.1\endcsname{8.11.1}
\expandafter\def\csname SGAsource@I.sec8.11.5\endcsname{8.11.5}
\expandafter\def\csname SGAsource@I.sec8.12\endcsname{8.12}
\expandafter\def\csname SGAsource@I.sec8.12.1\endcsname{8.12.1}
\expandafter\def\csname SGAsource@I.sec8.12.2.2\endcsname{8.12.2.2}
\expandafter\def\csname SGAsource@I.sec8.12.2.3\endcsname{8.12.2.3}
\expandafter\def\csname SGAsource@I.sec8.12.3\endcsname{8.12.3}
\expandafter\def\csname SGAsource@I.sec8.13\endcsname{8.13}
\expandafter\def\csname SGAsource@I.sec8.13.1\endcsname{8.13.1}
\expandafter\def\csname SGAsource@I.sec8.2\endcsname{8.2}
\expandafter\def\csname SGAsource@I.sec8.2.1\endcsname{8.2.1}
\expandafter\def\csname SGAsource@I.sec8.2.2\endcsname{8.2.2}
\expandafter\def\csname SGAsource@I.sec8.2.3\endcsname{8.2.3}
\expandafter\def\csname SGAsource@I.sec8.2.4\endcsname{8.2.4}
\expandafter\def\csname SGAsource@I.sec8.2.5\endcsname{8.2.5}
\expandafter\def\csname SGAsource@I.sec8.2.6\endcsname{8.2.6}
\expandafter\def\csname SGAsource@I.sec8.2.7\endcsname{8.2.7}
\expandafter\def\csname SGAsource@I.sec8.3\endcsname{8.3}
\expandafter\def\csname SGAsource@I.sec8.3.2\endcsname{8.3.2}
\expandafter\def\csname SGAsource@I.sec8.3.2.3\endcsname{8.3.2.3}
\expandafter\def\csname SGAsource@I.sec8.4\endcsname{8.4}
\expandafter\def\csname SGAsource@I.sec8.4.3\endcsname{8.4.3}
\expandafter\def\csname SGAsource@I.sec8.4.4\endcsname{8.4.4}
\expandafter\def\csname SGAsource@I.sec8.5\endcsname{8.5}
\expandafter\def\csname SGAsource@I.sec8.5.3\endcsname{8.5.3}
\expandafter\def\csname SGAsource@I.sec8.5.4\endcsname{8.5.4}
\expandafter\def\csname SGAsource@I.sec8.6\endcsname{8.6}
\expandafter\def\csname SGAsource@I.sec8.6.1\endcsname{8.6.1}
\expandafter\def\csname SGAsource@I.sec8.6.2\endcsname{8.6.2}
\expandafter\def\csname SGAsource@I.sec8.7\endcsname{8.7}
\expandafter\def\csname SGAsource@I.sec8.7.1\endcsname{8.7.1}
\expandafter\def\csname SGAsource@I.sec8.7.1.6\endcsname{8.7.1.6}
\expandafter\def\csname SGAsource@I.sec8.7.1.7\endcsname{8.7.1.7}
\expandafter\def\csname SGAsource@I.sec8.7.2\endcsname{8.7.2}
\expandafter\def\csname SGAsource@I.sec8.7.4\endcsname{8.7.4}
\expandafter\def\csname SGAsource@I.sec8.8\endcsname{8.8}
\expandafter\def\csname SGAsource@I.sec8.8.1\endcsname{8.8.1}
\expandafter\def\csname SGAsource@I.sec8.9\endcsname{8.9}
\expandafter\def\csname SGAsource@I.sec9\endcsname{9}
\expandafter\def\csname SGAsource@I.sec9.0\endcsname{9.0}
\expandafter\def\csname SGAsource@I.sec9.1.1\endcsname{9.1.1}
\expandafter\def\csname SGAsource@I.sec9.1.2\endcsname{9.1.2}
\expandafter\def\csname SGAsource@I.sec9.15.1\endcsname{9.15.1}
\expandafter\def\csname SGAsource@I.sec9.2.1\endcsname{9.2.1}
\expandafter\def\csname SGAsource@I.sec9.21\endcsname{9.21}
\expandafter\def\csname SGAsource@I.sec9.21.14\endcsname{9.21.14}
\expandafter\def\csname SGAsource@I.sec9.21.15\endcsname{9.21.15}
\expandafter\def\csname SGAsource@I.sec9.21.20\endcsname{9.21.20}
\expandafter\def\csname SGAsource@I.sec9.21.21\endcsname{9.21.21}
\expandafter\def\csname SGAsource@I.sec9.3.1\endcsname{9.3.1}
\expandafter\def\csname SGAsource@I.them8.3.3\endcsname{8.3.3}
\expandafter\def\csname SGAsource@I.them9.22\endcsname{9.22}
\expandafter\def\csname SGAsource@I.var7.5.2\endcsname{7.5.2}
\expandafter\def\csname SGAsource@II.cor2.3\endcsname{2.3}
\expandafter\def\csname SGAsource@II.cor2.4\endcsname{2.4}
\expandafter\def\csname SGAsource@II.cor4.1.1\endcsname{4.1.1}
\expandafter\def\csname SGAsource@II.cor4.11\endcsname{4.11}
\expandafter\def\csname SGAsource@II.cor4.12\endcsname{4.12}
\expandafter\def\csname SGAsource@II.cor4.4.4\endcsname{4.4.4}
\expandafter\def\csname SGAsource@II.cor4.6.1\endcsname{4.6.1}
\expandafter\def\csname SGAsource@II.cor4.6.2\endcsname{4.6.2}
\expandafter\def\csname SGAsource@II.cor4.6.3\endcsname{4.6.3}
\expandafter\def\csname SGAsource@II.cor4.6.4\endcsname{4.6.4}
\expandafter\def\csname SGAsource@II.cor6.6\endcsname{6.6}
\expandafter\def\csname SGAsource@II.def1.1\endcsname{1.1}
\expandafter\def\csname SGAsource@II.def1.2\endcsname{1.2}
\expandafter\def\csname SGAsource@II.def1.3\endcsname{1.3}
\expandafter\def\csname SGAsource@II.def2.1\endcsname{2.1}
\expandafter\def\csname SGAsource@II.def2.5\endcsname{2.5}
\expandafter\def\csname SGAsource@II.def3.0.1\endcsname{3.0.1}
\expandafter\def\csname SGAsource@II.def3.0.2\endcsname{3.0.2}
\expandafter\def\csname SGAsource@II.def3.5\endcsname{3.5}
\expandafter\def\csname SGAsource@II.def4.5\endcsname{4.5}
\expandafter\def\csname SGAsource@II.def5.2\endcsname{5.2}
\expandafter\def\csname SGAsource@II.def6.1\endcsname{6.1}
\expandafter\def\csname SGAsource@II.eq4.10.1\endcsname{4.10.1}
\expandafter\def\csname SGAsource@II.eq4.10.2\endcsname{4.10.2}
\expandafter\def\csname SGAsource@II.eq4.6.1\endcsname{4.6.1}
\expandafter\def\csname SGAsource@II.exem4.5.1\endcsname{4.5.1}
\expandafter\def\csname SGAsource@II.lem3.1\endcsname{3.1}
\expandafter\def\csname SGAsource@II.not6.8\endcsname{6.8}
\expandafter\def\csname SGAsource@II.prop1.4\endcsname{1.4}
\expandafter\def\csname SGAsource@II.prop2.2\endcsname{2.2}
\expandafter\def\csname SGAsource@II.prop2.6\endcsname{2.6}
\expandafter\def\csname SGAsource@II.prop3.0.4\endcsname{3.0.4}
\expandafter\def\csname SGAsource@II.prop3.2\endcsname{3.2}
\expandafter\def\csname SGAsource@II.prop3.6\endcsname{3.6}
\expandafter\def\csname SGAsource@II.prop4.10\endcsname{4.10}
\expandafter\def\csname SGAsource@II.prop4.2\endcsname{4.2}
\expandafter\def\csname SGAsource@II.prop4.3\endcsname{4.3}
\expandafter\def\csname SGAsource@II.prop4.6\endcsname{4.6}
\expandafter\def\csname SGAsource@II.prop4.7\endcsname{4.7}
\expandafter\def\csname SGAsource@II.prop4.8\endcsname{4.8}
\expandafter\def\csname SGAsource@II.prop5.1\endcsname{5.1}
\expandafter\def\csname SGAsource@II.prop5.3\endcsname{5.3}
\expandafter\def\csname SGAsource@II.prop5.4\endcsname{5.4}
\expandafter\def\csname SGAsource@II.prop6.3.1\endcsname{6.3.1}
\expandafter\def\csname SGAsource@II.prop6.4\endcsname{6.4}
\expandafter\def\csname SGAsource@II.prop6.5\endcsname{6.5}
\expandafter\def\csname SGAsource@II.prop6.7\endcsname{6.7}
\expandafter\def\csname SGAsource@II.rem2.5.1\endcsname{2.5.1}
\expandafter\def\csname SGAsource@II.rem2.7\endcsname{2.7}
\expandafter\def\csname SGAsource@II.rem3.3\endcsname{3.3}
\expandafter\def\csname SGAsource@II.rem4.3.1\endcsname{4.3.1}
\expandafter\def\csname SGAsource@II.rem5.4.1\endcsname{5.4.1}
\expandafter\def\csname SGAsource@II.rem5.4.2\endcsname{5.4.2}
\expandafter\def\csname SGAsource@II.rem6.2\endcsname{6.2}
\expandafter\def\csname SGAsource@II.rem6.9\endcsname{6.9}
\expandafter\def\csname SGAsource@II.sec1\endcsname{1}
\expandafter\def\csname SGAsource@II.sec2\endcsname{2}
\expandafter\def\csname SGAsource@II.sec3\endcsname{3}
\expandafter\def\csname SGAsource@II.sec4\endcsname{4}
\expandafter\def\csname SGAsource@II.sec5\endcsname{5}
\expandafter\def\csname SGAsource@II.sec6\endcsname{6}
\expandafter\def\csname SGAsource@II.subsec4.9\endcsname{4.9}
\expandafter\def\csname SGAsource@II.subsec6.0\endcsname{6.0}
\expandafter\def\csname SGAsource@II.subsubsec1.1.1\endcsname{1.1.1}
\expandafter\def\csname SGAsource@II.subsubsec1.1.2\endcsname{1.1.2}
\expandafter\def\csname SGAsource@II.subsubsec1.1.3\endcsname{1.1.3}
\expandafter\def\csname SGAsource@II.subsubsec1.1.4\endcsname{1.1.4}
\expandafter\def\csname SGAsource@II.subsubsec1.1.5\endcsname{1.1.5}
\expandafter\def\csname SGAsource@II.subsubsec1.1.6\endcsname{1.1.6}
\expandafter\def\csname SGAsource@II.subsubsec1.3.1\endcsname{1.3.1}
\expandafter\def\csname SGAsource@II.subsubsec3.0.3\endcsname{3.0.3}
\expandafter\def\csname SGAsource@II.subsubsec3.0.5\endcsname{3.0.5}
\expandafter\def\csname SGAsource@II.subsubsec4.4.0\endcsname{4.4.0}
\expandafter\def\csname SGAsource@II.subsubsec5.3.1\endcsname{5.3.1}
\expandafter\def\csname SGAsource@II.subsubsec5.5.0\endcsname{5.5.0}
\expandafter\def\csname SGAsource@II.subsubsec6.3.0\endcsname{6.3.0}
\expandafter\def\csname SGAsource@II.subsubsec6.3.2\endcsname{6.3.2}
\expandafter\def\csname SGAsource@II.subsubsec6.3.3\endcsname{6.3.3}
\expandafter\def\csname SGAsource@II.theo3.4\endcsname{3.4}
\expandafter\def\csname SGAsource@II.theo4.1\endcsname{4.1}
\expandafter\def\csname SGAsource@II.theo4.4\endcsname{4.4}
\expandafter\def\csname SGAsource@II.theo5.5\endcsname{5.5}
\expandafter\def\csname SGAsource@III.cor3.3\endcsname{3.3}
\expandafter\def\csname SGAsource@III.cor3.4\endcsname{3.4}
\expandafter\def\csname SGAsource@III.def1.1\endcsname{1.1}
\expandafter\def\csname SGAsource@III.def2.1\endcsname{2.1}
\expandafter\def\csname SGAsource@III.eq1.5.1\endcsname{1.5.1}
\expandafter\def\csname SGAsource@III.eq2.3.1\endcsname{2.3.1}
\expandafter\def\csname SGAsource@III.eq2.3.2\endcsname{2.3.2}
\expandafter\def\csname SGAsource@III.exem1.9.1\endcsname{1.9.1}
\expandafter\def\csname SGAsource@III.exem1.9.2\endcsname{1.9.2}
\expandafter\def\csname SGAsource@III.exem1.9.3\endcsname{1.9.3}
\expandafter\def\csname SGAsource@III.exem2.7\endcsname{2.7}
\expandafter\def\csname SGAsource@III.lem3.6.1\endcsname{3.6.1}
\expandafter\def\csname SGAsource@III.not1.8\endcsname{1.8}
\expandafter\def\csname SGAsource@III.prop1.2\endcsname{1.2}
\expandafter\def\csname SGAsource@III.prop1.3\endcsname{1.3}
\expandafter\def\csname SGAsource@III.prop1.5\endcsname{1.5}
\expandafter\def\csname SGAsource@III.prop1.6\endcsname{1.6}
\expandafter\def\csname SGAsource@III.prop1.7\endcsname{1.7}
\expandafter\def\csname SGAsource@III.prop2.2\endcsname{2.2}
\expandafter\def\csname SGAsource@III.prop2.3\endcsname{2.3}
\expandafter\def\csname SGAsource@III.prop2.5\endcsname{2.5}
\expandafter\def\csname SGAsource@III.prop2.6\endcsname{2.6}
\expandafter\def\csname SGAsource@III.prop3.2\endcsname{3.2}
\expandafter\def\csname SGAsource@III.prop3.5\endcsname{3.5}
\expandafter\def\csname SGAsource@III.prop3.6\endcsname{3.6}
\expandafter\def\csname SGAsource@III.prop3.7\endcsname{3.7}
\expandafter\def\csname SGAsource@III.prop5.2\endcsname{5.2}
\expandafter\def\csname SGAsource@III.prop5.4\endcsname{5.4}
\expandafter\def\csname SGAsource@III.prop5.5\endcsname{5.5}
\expandafter\def\csname SGAsource@III.rem1.4\endcsname{1.4}
\expandafter\def\csname SGAsource@III.sec1\endcsname{1}
\expandafter\def\csname SGAsource@III.sec2\endcsname{2}
\expandafter\def\csname SGAsource@III.sec3\endcsname{3}
\expandafter\def\csname SGAsource@III.sec4\endcsname{4}
\expandafter\def\csname SGAsource@III.sec5\endcsname{5}
\expandafter\def\csname SGAsource@III.subsec1.11\endcsname{1.11}
\expandafter\def\csname SGAsource@III.subsec2.4\endcsname{2.4}
\expandafter\def\csname SGAsource@III.subsec3.1\endcsname{3.1}
\expandafter\def\csname SGAsource@III.subsec4.2\endcsname{4.2}
\expandafter\def\csname SGAsource@III.subsec4.3\endcsname{4.3}
\expandafter\def\csname SGAsource@III.subsec4.4\endcsname{4.4}
\expandafter\def\csname SGAsource@III.subsec5.1\endcsname{5.1}
\expandafter\def\csname SGAsource@III.subsec5.3\endcsname{3}
\expandafter\def\csname SGAsource@III.subsubsec1.2.1\endcsname{1.2.1}
\expandafter\def\csname SGAsource@III.subsubsec4.1.1\endcsname{4.1.1}
\expandafter\def\csname SGAsource@III.subsubsec4.1.2\endcsname{4.1.2}
\expandafter\def\csname SGAsource@III.theo4.1\endcsname{4.1}
\expandafter\def\csname SGAsource@IV.cor1.2.1\endcsname{1.2.1}
\expandafter\def\csname SGAsource@IV.cor1.5\endcsname{1.5}
\expandafter\def\csname SGAsource@IV.cor1.6\endcsname{1.6}
\expandafter\def\csname SGAsource@IV.cor1.7\endcsname{1.7}
\expandafter\def\csname SGAsource@IV.cor10.4\endcsname{10.4}
\expandafter\def\csname SGAsource@IV.cor10.7\endcsname{10.7}
\expandafter\def\csname SGAsource@IV.cor10.8\endcsname{10.8}
\expandafter\def\csname SGAsource@IV.cor12.12\endcsname{12.12}
\expandafter\def\csname SGAsource@IV.cor12.13\endcsname{12.13}
\expandafter\def\csname SGAsource@IV.cor12.4\endcsname{12.4}
\expandafter\def\csname SGAsource@IV.cor12.5\endcsname{12.5}
\expandafter\def\csname SGAsource@IV.cor12.6\endcsname{12.6}
\expandafter\def\csname SGAsource@IV.cor13.5\endcsname{13.5}
\expandafter\def\csname SGAsource@IV.cor13.6\endcsname{13.6}
\expandafter\def\csname SGAsource@IV.cor4.2.4\endcsname{4.2.4}
\expandafter\def\csname SGAsource@IV.cor6.5.1\endcsname{6.5.1}
\expandafter\def\csname SGAsource@IV.cor6.7.3\endcsname{6.7.3}
\expandafter\def\csname SGAsource@IV.cor9.1.5\endcsname{9.1.5}
\expandafter\def\csname SGAsource@IV.cor9.4.3\endcsname{9.4.3}
\expandafter\def\csname SGAsource@IV.cor9.4.4\endcsname{9.4.4}
\expandafter\def\csname SGAsource@IV.cor9.4.5\endcsname{9.4.5}
\expandafter\def\csname SGAsource@IV.cor9.4.6\endcsname{9.4.6}
\expandafter\def\csname SGAsource@IV.cor9.7.3\endcsname{9.7.3}
\expandafter\def\csname SGAsource@IV.def1.1\endcsname{1.1}
\expandafter\def\csname SGAsource@IV.def13.1\endcsname{13.1}
\expandafter\def\csname SGAsource@IV.def13.3\endcsname{13.3}
\expandafter\def\csname SGAsource@IV.def3.1\endcsname{3.1}
\expandafter\def\csname SGAsource@IV.def6.1\endcsname{6.1}
\expandafter\def\csname SGAsource@IV.def6.2\endcsname{6.2}
\expandafter\def\csname SGAsource@IV.def6.4.1\endcsname{6.4.1}
\expandafter\def\csname SGAsource@IV.def8.3\endcsname{8.3}
\expandafter\def\csname SGAsource@IV.def9.1.1\endcsname{9.1.1}
\expandafter\def\csname SGAsource@IV.def9.5.5\endcsname{9.5.5}
\expandafter\def\csname SGAsource@IV.eq10.2.1\endcsname{10.2.1}
\expandafter\def\csname SGAsource@IV.eq10.3.1\endcsname{10.3.1}
\expandafter\def\csname SGAsource@IV.eq11.2.2.1\endcsname{11.2.2.1}
\expandafter\def\csname SGAsource@IV.eq12.1.1\endcsname{12.1.1}
\expandafter\def\csname SGAsource@IV.eq12.1.2\endcsname{12.1.2}
\expandafter\def\csname SGAsource@IV.eq12.14.1\endcsname{12.14.1}
\expandafter\def\csname SGAsource@IV.eq12.8.1\endcsname{12.8.1}
\expandafter\def\csname SGAsource@IV.eq13.4.3\endcsname{13.4.3}
\expandafter\def\csname SGAsource@IV.eq14.10.1\endcsname{14.10.1}
\expandafter\def\csname SGAsource@IV.eq14.2.1\endcsname{14.2.1}
\expandafter\def\csname SGAsource@IV.eq14.4.1\endcsname{14.4.1}
\expandafter\def\csname SGAsource@IV.eq14.4.2\endcsname{14.4.2}
\expandafter\def\csname SGAsource@IV.eq4.1.1.1\endcsname{4.1.1.1}
\expandafter\def\csname SGAsource@IV.eq4.10.1\endcsname{4.10.1}
\expandafter\def\csname SGAsource@IV.eq4.10.1.1\endcsname{4.10.1.1}
\expandafter\def\csname SGAsource@IV.eq4.10.2.1\endcsname{4.10.2.1}
\expandafter\def\csname SGAsource@IV.eq4.10.3.1\endcsname{4.10.3.1}
\expandafter\def\csname SGAsource@IV.eq4.2.1.1\endcsname{4.2.1.1}
\expandafter\def\csname SGAsource@IV.eq4.3.2\endcsname{4.3.2}
\expandafter\def\csname SGAsource@IV.eq4.3.3\endcsname{4.3.3}
\expandafter\def\csname SGAsource@IV.eq4.4.1\endcsname{4.4.1}
\expandafter\def\csname SGAsource@IV.eq4.4.2\endcsname{4.4.2}
\expandafter\def\csname SGAsource@IV.eq4.6.2.1\endcsname{4.6.2.1}
\expandafter\def\csname SGAsource@IV.eq4.6.2.2\endcsname{4.6.2.2}
\expandafter\def\csname SGAsource@IV.eq4.6.3.1\endcsname{4.6.3.1}
\expandafter\def\csname SGAsource@IV.eq4.6.3.2\endcsname{4.6.3.2}
\expandafter\def\csname SGAsource@IV.eq4.7.1\endcsname{4.7.1}
\expandafter\def\csname SGAsource@IV.eq4.7.2\endcsname{4.7.2}
\expandafter\def\csname SGAsource@IV.eq4.7.3\endcsname{4.7.3}
\expandafter\def\csname SGAsource@IV.eq4.9.1.1\endcsname{4.9.1.1}
\expandafter\def\csname SGAsource@IV.eq5.10.1\endcsname{5.10.1}
\expandafter\def\csname SGAsource@IV.eq5.10.2\endcsname{5.10.2}
\expandafter\def\csname SGAsource@IV.eq5.10.3\endcsname{5.10.3}
\expandafter\def\csname SGAsource@IV.eq5.11.1.1\endcsname{5.11.1.1}
\expandafter\def\csname SGAsource@IV.eq5.11.1.2\endcsname{5.11.1.2}
\expandafter\def\csname SGAsource@IV.eq5.12.1\endcsname{5.12.1}
\expandafter\def\csname SGAsource@IV.eq5.2.1\endcsname{5.2.1}
\expandafter\def\csname SGAsource@IV.eq5.2.2\endcsname{5.2.2}
\expandafter\def\csname SGAsource@IV.eq5.4.1\endcsname{5.4.1}
\expandafter\def\csname SGAsource@IV.eq5.5.1\endcsname{5.5.1}
\expandafter\def\csname SGAsource@IV.eq5.5.2\endcsname{5.5.2}
\expandafter\def\csname SGAsource@IV.eq5.5.3\endcsname{5.5.3}
\expandafter\def\csname SGAsource@IV.eq5.8.1\endcsname{5.8.1}
\expandafter\def\csname SGAsource@IV.eq5.8.2\endcsname{5.8.2}
\expandafter\def\csname SGAsource@IV.eq5.8.3.1\endcsname{5.8.3.1}
\expandafter\def\csname SGAsource@IV.eq5.8.3.2\endcsname{5.8.3.2}
\expandafter\def\csname SGAsource@IV.eq6.10.1\endcsname{6.10.1}
\expandafter\def\csname SGAsource@IV.eq6.2.1.1\endcsname{6.2.1.1}
\expandafter\def\csname SGAsource@IV.eq6.3.1\endcsname{6.3.1}
\expandafter\def\csname SGAsource@IV.eq6.3.2\endcsname{6.3.2}
\expandafter\def\csname SGAsource@IV.eq6.3.3\endcsname{6.3.3}
\expandafter\def\csname SGAsource@IV.eq6.7.1\endcsname{6.7.1}
\expandafter\def\csname SGAsource@IV.eq6.7.2\endcsname{6.7.2}
\expandafter\def\csname SGAsource@IV.eq6.8.1\endcsname{6.8.1}
\expandafter\def\csname SGAsource@IV.eq6.8.3\endcsname{6.8.3}
\expandafter\def\csname SGAsource@IV.eq6.8.4\endcsname{6.8.4}
\expandafter\def\csname SGAsource@IV.eq6.8.5.1\endcsname{6.8.5.1}
\expandafter\def\csname SGAsource@IV.eq6.8.6.1\endcsname{6.8.6.1}
\expandafter\def\csname SGAsource@IV.eq6.8.7.1\endcsname{6.8.7.1}
\expandafter\def\csname SGAsource@IV.eq7.1.1\endcsname{7.1.1}
\expandafter\def\csname SGAsource@IV.eq7.1.2\endcsname{7.1.2}
\expandafter\def\csname SGAsource@IV.eq7.1.3\endcsname{7.1.3}
\expandafter\def\csname SGAsource@IV.eq7.1.4\endcsname{7.1.4}
\expandafter\def\csname SGAsource@IV.eq7.2.1\endcsname{7.2.1}
\expandafter\def\csname SGAsource@IV.eq7.2.2\endcsname{7.2.2}
\expandafter\def\csname SGAsource@IV.eq7.2.3\endcsname{7.2.3}
\expandafter\def\csname SGAsource@IV.eq7.2.4.1\endcsname{7.2.4.1}
\expandafter\def\csname SGAsource@IV.eq7.2.6.1\endcsname{7.2.6.1}
\expandafter\def\csname SGAsource@IV.eq7.2.6.2\endcsname{7.2.6.2}
\expandafter\def\csname SGAsource@IV.eq8.7.1\endcsname{8.7.1}
\expandafter\def\csname SGAsource@IV.eq8.7.2\endcsname{8.7.2}
\expandafter\def\csname SGAsource@IV.eq8.7.3\endcsname{8.7.3}
\expandafter\def\csname SGAsource@IV.eq8.7.4\endcsname{8.7.4}
\expandafter\def\csname SGAsource@IV.eq8.7.5\endcsname{8.7.5}
\expandafter\def\csname SGAsource@IV.eq8.7.6\endcsname{8.7.6}
\expandafter\def\csname SGAsource@IV.eq9.1.12.2\endcsname{9.1.12.2}
\expandafter\def\csname SGAsource@IV.eq9.1.14.1\endcsname{9.1.14.1}
\expandafter\def\csname SGAsource@IV.eq9.1.4.1\endcsname{9.1.4.1}
\expandafter\def\csname SGAsource@IV.eq9.1.5.1\endcsname{9.1.5.1}
\expandafter\def\csname SGAsource@IV.eq9.1.6.1\endcsname{9.1.6.1}
\expandafter\def\csname SGAsource@IV.eq9.1.7.1\endcsname{9.1.7.1}
\expandafter\def\csname SGAsource@IV.eq9.1.7.3\endcsname{9.1.7.3}
\expandafter\def\csname SGAsource@IV.eq9.1.8.1\endcsname{9.1.8.1}
\expandafter\def\csname SGAsource@IV.eq9.1.8.2\endcsname{9.1.8.2}
\expandafter\def\csname SGAsource@IV.eq9.2.2\endcsname{9.2.2}
\expandafter\def\csname SGAsource@IV.eq9.2.6.1\endcsname{9.2.6.1}
\expandafter\def\csname SGAsource@IV.eq9.3.2.1\endcsname{9.3.2.1}
\expandafter\def\csname SGAsource@IV.eq9.3.2.2\endcsname{9.3.2.2}
\expandafter\def\csname SGAsource@IV.eq9.3.2.3\endcsname{9.3.2.3}
\expandafter\def\csname SGAsource@IV.eq9.3.6.1\endcsname{9.3.6.1}
\expandafter\def\csname SGAsource@IV.eq9.3.6.2\endcsname{9.3.6.2}
\expandafter\def\csname SGAsource@IV.eq9.3.6.3\endcsname{9.3.6.3}
\expandafter\def\csname SGAsource@IV.eq9.4.3.1\endcsname{9.4.3.1}
\expandafter\def\csname SGAsource@IV.eq9.5.10.1\endcsname{9.5.10.1}
\expandafter\def\csname SGAsource@IV.eq9.5.10.2\endcsname{9.5.10.2}
\expandafter\def\csname SGAsource@IV.eq9.5.10.3\endcsname{9.5.10.3}
\expandafter\def\csname SGAsource@IV.eq9.5.3.1\endcsname{9.5.3.1}
\expandafter\def\csname SGAsource@IV.eq9.6.1\endcsname{9.6.1}
\expandafter\def\csname SGAsource@IV.eq9.6.2\endcsname{9.6.2}
\expandafter\def\csname SGAsource@IV.eq9.6.3\endcsname{9.6.3}
\expandafter\def\csname SGAsource@IV.eq9.6.4\endcsname{9.6.4}
\expandafter\def\csname SGAsource@IV.eq9.7.2.1\endcsname{9.7.2.1}
\expandafter\def\csname SGAsource@IV.eq9.7.2.2\endcsname{9.7.2.2}
\expandafter\def\csname SGAsource@IV.eq9.7.2.3\endcsname{9.7.2.3}
\expandafter\def\csname SGAsource@IV.eq9.7.2.4.1\endcsname{9.7.2.4.1}
\expandafter\def\csname SGAsource@IV.eq9.8.2.1\endcsname{9.8.2.1}
\expandafter\def\csname SGAsource@IV.exer11.3.5\endcsname{11.3.5}
\expandafter\def\csname SGAsource@IV.exer13.9\endcsname{13.9}
\expandafter\def\csname SGAsource@IV.exer2.7.5\endcsname{2.7.5}
\expandafter\def\csname SGAsource@IV.exer4.10.6\endcsname{4.10.6}
\expandafter\def\csname SGAsource@IV.exer5.14\endcsname{5.14}
\expandafter\def\csname SGAsource@IV.exer5.9\endcsname{5.9}
\expandafter\def\csname SGAsource@IV.exer6.10\endcsname{6.10}
\expandafter\def\csname SGAsource@IV.exer6.5.2\endcsname{6.5.2}
\expandafter\def\csname SGAsource@IV.exer6.9\endcsname{6.9}
\expandafter\def\csname SGAsource@IV.exer7.1.10\endcsname{7.1.10}
\expandafter\def\csname SGAsource@IV.exer7.1.11\endcsname{7.1.11}
\expandafter\def\csname SGAsource@IV.exer7.1.9\endcsname{7.1.9}
\expandafter\def\csname SGAsource@IV.exer7.2.6\endcsname{7.2.6}
\expandafter\def\csname SGAsource@IV.exer7.5\endcsname{7.5}
\expandafter\def\csname SGAsource@IV.exer7.6\endcsname{7.6}
\expandafter\def\csname SGAsource@IV.exer7.7\endcsname{7.7}
\expandafter\def\csname SGAsource@IV.exer7.8\endcsname{7.8}
\expandafter\def\csname SGAsource@IV.exer8.6\endcsname{8.6}
\expandafter\def\csname SGAsource@IV.exer8.7\endcsname{8.7}
\expandafter\def\csname SGAsource@IV.exer8.8\endcsname{8.8}
\expandafter\def\csname SGAsource@IV.exer9.1.10\endcsname{9.1.10}
\expandafter\def\csname SGAsource@IV.exer9.1.11\endcsname{9.1.11}
\expandafter\def\csname SGAsource@IV.exer9.1.12\endcsname{9.1.12}
\expandafter\def\csname SGAsource@IV.exer9.1.13\endcsname{9.1.13}
\expandafter\def\csname SGAsource@IV.exer9.1.6\endcsname{9.1.6}
\expandafter\def\csname SGAsource@IV.exer9.1.7\endcsname{9.1.7}
\expandafter\def\csname SGAsource@IV.exer9.1.7.2\endcsname{9.1.7.2}
\expandafter\def\csname SGAsource@IV.exer9.1.8\endcsname{9.1.8}
\expandafter\def\csname SGAsource@IV.exer9.1.9\endcsname{9.1.9}
\expandafter\def\csname SGAsource@IV.exer9.2.7\endcsname{9.2.7}
\expandafter\def\csname SGAsource@IV.exer9.4.10\endcsname{9.4.10}
\expandafter\def\csname SGAsource@IV.exer9.4.7\endcsname{9.4.7}
\expandafter\def\csname SGAsource@IV.exer9.4.8\endcsname{9.4.8}
\expandafter\def\csname SGAsource@IV.exer9.4.9\endcsname{9.4.9}
\expandafter\def\csname SGAsource@IV.exer9.5.10\endcsname{9.5.10}
\expandafter\def\csname SGAsource@IV.exer9.5.11\endcsname{9.5.11}
\expandafter\def\csname SGAsource@IV.exer9.5.12\endcsname{9.5.12}
\expandafter\def\csname SGAsource@IV.exer9.5.9\endcsname{9.5.9}
\expandafter\def\csname SGAsource@IV.exer9.6.5\endcsname{9.6.5}
\expandafter\def\csname SGAsource@IV.exer9.7.5\endcsname{9.7.5}
\expandafter\def\csname SGAsource@IV.exer9.8.1\endcsname{9.8.1}
\expandafter\def\csname SGAsource@IV.exer9.8.2\endcsname{9.8.2}
\expandafter\def\csname SGAsource@IV.exer9.8.3\endcsname{9.8.3}
\expandafter\def\csname SGAsource@IV.lem10.4.2\endcsname{10.4.2}
\expandafter\def\csname SGAsource@IV.lem5.11.2\endcsname{5.11.2}
\expandafter\def\csname SGAsource@IV.lem9.7.2.4\endcsname{9.7.2.4}
\expandafter\def\csname SGAsource@IV.par1.2.4.1\endcsname{1.2.4.1}
\expandafter\def\csname SGAsource@IV.par1.2.4.2\endcsname{1.2.4.2}
\expandafter\def\csname SGAsource@IV.par1.2.4.3\endcsname{1.2.4.3}
\expandafter\def\csname SGAsource@IV.par1.2.4.4\endcsname{1.2.4.4}
\expandafter\def\csname SGAsource@IV.par9.5.4.8\endcsname{9.5.4.8}
\expandafter\def\csname SGAsource@IV.par9.5.8.1\endcsname{9.5.8.1}
\expandafter\def\csname SGAsource@IV.par9.5.8.2\endcsname{9.5.8.2}
\expandafter\def\csname SGAsource@IV.par9.5.8.3\endcsname{9.5.8.3}
\expandafter\def\csname SGAsource@IV.par9.5.8.4\endcsname{9.5.8.4}
\expandafter\def\csname SGAsource@IV.par9.5.8.5\endcsname{9.5.8.5}
\expandafter\def\csname SGAsource@IV.par9.5.8.6\endcsname{9.5.8.6}
\expandafter\def\csname SGAsource@IV.par9.5.8.7\endcsname{9.5.8.7}
\expandafter\def\csname SGAsource@IV.par9.5.8.8\endcsname{9.5.8.8}
\expandafter\def\csname SGAsource@IV.prop1.4\endcsname{1.4}
\expandafter\def\csname SGAsource@IV.prop10.1\endcsname{10.1}
\expandafter\def\csname SGAsource@IV.prop10.3\endcsname{10.3}
\expandafter\def\csname SGAsource@IV.prop10.6\endcsname{10.6}
\expandafter\def\csname SGAsource@IV.prop11.1.6\endcsname{11.1.6}
\expandafter\def\csname SGAsource@IV.prop11.3.1\endcsname{11.3.1}
\expandafter\def\csname SGAsource@IV.prop11.3.3\endcsname{11.3.3}
\expandafter\def\csname SGAsource@IV.prop12.1\endcsname{12.1}
\expandafter\def\csname SGAsource@IV.prop12.10\endcsname{12.10}
\expandafter\def\csname SGAsource@IV.prop12.11\endcsname{12.11}
\expandafter\def\csname SGAsource@IV.prop12.14\endcsname{12.14}
\expandafter\def\csname SGAsource@IV.prop12.3\endcsname{12.3}
\expandafter\def\csname SGAsource@IV.prop12.7\endcsname{12.7}
\expandafter\def\csname SGAsource@IV.prop13.4\endcsname{13.4}
\expandafter\def\csname SGAsource@IV.prop14.10\endcsname{14.10}
\expandafter\def\csname SGAsource@IV.prop14.5\endcsname{14.5}
\expandafter\def\csname SGAsource@IV.prop14.6\endcsname{14.6}
\expandafter\def\csname SGAsource@IV.prop14.8\endcsname{14.8}
\expandafter\def\csname SGAsource@IV.prop4.9.4\endcsname{4.9.4}
\expandafter\def\csname SGAsource@IV.prop5.11\endcsname{5.11}
\expandafter\def\csname SGAsource@IV.prop5.12\endcsname{5.12}
\expandafter\def\csname SGAsource@IV.prop6.5\endcsname{6.5}
\expandafter\def\csname SGAsource@IV.prop6.5.3\endcsname{6.5.3}
\expandafter\def\csname SGAsource@IV.prop9.1.3\endcsname{9.1.3}
\expandafter\def\csname SGAsource@IV.prop9.1.4\endcsname{9.1.4}
\expandafter\def\csname SGAsource@IV.prop9.2.4\endcsname{9.2.4}
\expandafter\def\csname SGAsource@IV.prop9.3.3\endcsname{9.3.3}
\expandafter\def\csname SGAsource@IV.prop9.3.4\endcsname{9.3.4}
\expandafter\def\csname SGAsource@IV.prop9.4.1\endcsname{9.4.1}
\expandafter\def\csname SGAsource@IV.prop9.4.2\endcsname{9.4.2}
\expandafter\def\csname SGAsource@IV.prop9.7.2\endcsname{9.7.2}
\expandafter\def\csname SGAsource@IV.rem1.3\endcsname{1.3}
\expandafter\def\csname SGAsource@IV.rem1.8\endcsname{1.8}
\expandafter\def\csname SGAsource@IV.rem11.3.2\endcsname{11.3.2}
\expandafter\def\csname SGAsource@IV.rem11.3.4\endcsname{11.3.4}
\expandafter\def\csname SGAsource@IV.rem4.10.5\endcsname{4.10.5}
\expandafter\def\csname SGAsource@IV.rem4.6.4\endcsname{4.6.4}
\expandafter\def\csname SGAsource@IV.rem4.7.4\endcsname{4.7.4}
\expandafter\def\csname SGAsource@IV.rem5.13\endcsname{5.13}
\expandafter\def\csname SGAsource@IV.rem6.7.4\endcsname{6.7.4}
\expandafter\def\csname SGAsource@IV.rem7.2.4\endcsname{7.2.4}
\expandafter\def\csname SGAsource@IV.rem8.1.1\endcsname{8.1.1}
\expandafter\def\csname SGAsource@IV.rem8.2.1\endcsname{8.2.1}
\expandafter\def\csname SGAsource@IV.rem8.5.3\endcsname{8.5.3}
\expandafter\def\csname SGAsource@IV.rem9.1.2\endcsname{9.1.2}
\expandafter\def\csname SGAsource@IV.rem9.2.3\endcsname{9.2.3}
\expandafter\def\csname SGAsource@IV.rem9.2.4.1\endcsname{9.2.4.1}
\expandafter\def\csname SGAsource@IV.rem9.7.4\endcsname{9.7.4}
\expandafter\def\csname SGAsource@IV.sec0\endcsname{0}
\expandafter\def\csname SGAsource@IV.sec1\endcsname{1}
\expandafter\def\csname SGAsource@IV.sec10\endcsname{10}
\expandafter\def\csname SGAsource@IV.sec11\endcsname{11}
\expandafter\def\csname SGAsource@IV.sec12\endcsname{12}
\expandafter\def\csname SGAsource@IV.sec13\endcsname{13}
\expandafter\def\csname SGAsource@IV.sec14\endcsname{14}
\expandafter\def\csname SGAsource@IV.sec2\endcsname{2}
\expandafter\def\csname SGAsource@IV.sec3\endcsname{3}
\expandafter\def\csname SGAsource@IV.sec4\endcsname{4}
\expandafter\def\csname SGAsource@IV.sec5\endcsname{5}
\expandafter\def\csname SGAsource@IV.sec6\endcsname{6}
\expandafter\def\csname SGAsource@IV.sec7\endcsname{7}
\expandafter\def\csname SGAsource@IV.sec8\endcsname{8}
\expandafter\def\csname SGAsource@IV.sec9\endcsname{9}
\expandafter\def\csname SGAsource@IV.subsec0.1\endcsname{0.1}
\expandafter\def\csname SGAsource@IV.subsec0.2\endcsname{0.2}
\expandafter\def\csname SGAsource@IV.subsec0.3\endcsname{0.3}
\expandafter\def\csname SGAsource@IV.subsec0.4\endcsname{0.4}
\expandafter\def\csname SGAsource@IV.subsec0.5\endcsname{0.5}
\expandafter\def\csname SGAsource@IV.subsec10.2\endcsname{10.2}
\expandafter\def\csname SGAsource@IV.subsec10.5\endcsname{10.5}
\expandafter\def\csname SGAsource@IV.subsec12.2\endcsname{12.2}
\expandafter\def\csname SGAsource@IV.subsec12.8\endcsname{12.8}
\expandafter\def\csname SGAsource@IV.subsec12.9\endcsname{12.9}
\expandafter\def\csname SGAsource@IV.subsec13.2\endcsname{13.2}
\expandafter\def\csname SGAsource@IV.subsec13.7\endcsname{13.7}
\expandafter\def\csname SGAsource@IV.subsec13.8\endcsname{13.8}
\expandafter\def\csname SGAsource@IV.subsec14.1\endcsname{14.1}
\expandafter\def\csname SGAsource@IV.subsec14.2\endcsname{14.2}
\expandafter\def\csname SGAsource@IV.subsec14.3\endcsname{14.3}
\expandafter\def\csname SGAsource@IV.subsec14.4\endcsname{14.4}
\expandafter\def\csname SGAsource@IV.subsec14.7\endcsname{14.7}
\expandafter\def\csname SGAsource@IV.subsec14.9\endcsname{14.9}
\expandafter\def\csname SGAsource@IV.subsec2.0\endcsname{2.0}
\expandafter\def\csname SGAsource@IV.subsec2.1\endcsname{2.1}
\expandafter\def\csname SGAsource@IV.subsec2.2\endcsname{2.2}
\expandafter\def\csname SGAsource@IV.subsec2.3\endcsname{2.3}
\expandafter\def\csname SGAsource@IV.subsec2.4\endcsname{2.4}
\expandafter\def\csname SGAsource@IV.subsec2.5\endcsname{2.5}
\expandafter\def\csname SGAsource@IV.subsec2.6\endcsname{2.6}
\expandafter\def\csname SGAsource@IV.subsec2.7\endcsname{2.7}
\expandafter\def\csname SGAsource@IV.subsec2.8\endcsname{2.8}
\expandafter\def\csname SGAsource@IV.subsec3.2\endcsname{3.2}
\expandafter\def\csname SGAsource@IV.subsec3.4\endcsname{3.4}
\expandafter\def\csname SGAsource@IV.subsec4.1\endcsname{4.1}
\expandafter\def\csname SGAsource@IV.subsec4.10\endcsname{4.10}
\expandafter\def\csname SGAsource@IV.subsec4.2\endcsname{4.2}
\expandafter\def\csname SGAsource@IV.subsec4.3\endcsname{4.3}
\expandafter\def\csname SGAsource@IV.subsec4.4\endcsname{4.4}
\expandafter\def\csname SGAsource@IV.subsec4.5\endcsname{4.5}
\expandafter\def\csname SGAsource@IV.subsec4.6\endcsname{4.6}
\expandafter\def\csname SGAsource@IV.subsec4.7\endcsname{4.7}
\expandafter\def\csname SGAsource@IV.subsec4.8\endcsname{4.8}
\expandafter\def\csname SGAsource@IV.subsec4.9\endcsname{4.9}
\expandafter\def\csname SGAsource@IV.subsec5.1\endcsname{5.1}
\expandafter\def\csname SGAsource@IV.subsec5.10\endcsname{5.10}
\expandafter\def\csname SGAsource@IV.subsec5.2\endcsname{5.2}
\expandafter\def\csname SGAsource@IV.subsec5.3\endcsname{5.3}
\expandafter\def\csname SGAsource@IV.subsec5.4\endcsname{5.4}
\expandafter\def\csname SGAsource@IV.subsec5.5\endcsname{5.5}
\expandafter\def\csname SGAsource@IV.subsec5.6\endcsname{5.6}
\expandafter\def\csname SGAsource@IV.subsec5.7\endcsname{5.7}
\expandafter\def\csname SGAsource@IV.subsec5.8\endcsname{5.8}
\expandafter\def\csname SGAsource@IV.subsec6.0\endcsname{6.0}
\expandafter\def\csname SGAsource@IV.subsec6.3\endcsname{6.3}
\expandafter\def\csname SGAsource@IV.subsec6.6\endcsname{6.6}
\expandafter\def\csname SGAsource@IV.subsec6.7\endcsname{6.7}
\expandafter\def\csname SGAsource@IV.subsec6.8\endcsname{6.8}
\expandafter\def\csname SGAsource@IV.subsec7.1\endcsname{7.1}
\expandafter\def\csname SGAsource@IV.subsec7.2\endcsname{7.2}
\expandafter\def\csname SGAsource@IV.subsec7.3\endcsname{7.3}
\expandafter\def\csname SGAsource@IV.subsec7.4\endcsname{7.4}
\expandafter\def\csname SGAsource@IV.subsec8.1\endcsname{8.1}
\expandafter\def\csname SGAsource@IV.subsec8.2\endcsname{8.2}
\expandafter\def\csname SGAsource@IV.subsec8.4\endcsname{8.4}
\expandafter\def\csname SGAsource@IV.subsec8.5\endcsname{8.5}
\expandafter\def\csname SGAsource@IV.subsec9.1\endcsname{9.1}
\expandafter\def\csname SGAsource@IV.subsec9.2\endcsname{9.2}
\expandafter\def\csname SGAsource@IV.subsec9.3\endcsname{9.3}
\expandafter\def\csname SGAsource@IV.subsec9.4\endcsname{9.4}
\expandafter\def\csname SGAsource@IV.subsec9.5\endcsname{9.5}
\expandafter\def\csname SGAsource@IV.subsec9.6\endcsname{9.6}
\expandafter\def\csname SGAsource@IV.subsec9.7\endcsname{9.7}
\expandafter\def\csname SGAsource@IV.subsec9.7.1\endcsname{9.7.1}
\expandafter\def\csname SGAsource@IV.subsec9.8\endcsname{9.8}
\expandafter\def\csname SGAsource@IV.subsubsec1.1.1\endcsname{1.1.1}
\expandafter\def\csname SGAsource@IV.subsubsec1.1.2\endcsname{1.1.2}
\expandafter\def\csname SGAsource@IV.subsubsec1.2.3\endcsname{1.2.3}
\expandafter\def\csname SGAsource@IV.subsubsec1.2.4\endcsname{1.2.4}
\expandafter\def\csname SGAsource@IV.subsubsec10.4.1\endcsname{10.4.1}
\expandafter\def\csname SGAsource@IV.subsubsec10.4.3\endcsname{10.4.3}
\expandafter\def\csname SGAsource@IV.subsubsec10.5.1\endcsname{10.5.1}
\expandafter\def\csname SGAsource@IV.subsubsec10.5.2\endcsname{10.5.2}
\expandafter\def\csname SGAsource@IV.subsubsec10.5.3\endcsname{10.5.3}
\expandafter\def\csname SGAsource@IV.subsubsec10.5.4\endcsname{10.5.4}
\expandafter\def\csname SGAsource@IV.subsubsec10.6.1\endcsname{10.6.1}
\expandafter\def\csname SGAsource@IV.subsubsec11.1.1\endcsname{11.1.1}
\expandafter\def\csname SGAsource@IV.subsubsec11.1.2\endcsname{11.1.2}
\expandafter\def\csname SGAsource@IV.subsubsec11.1.3\endcsname{11.1.3}
\expandafter\def\csname SGAsource@IV.subsubsec11.1.4\endcsname{11.1.4}
\expandafter\def\csname SGAsource@IV.subsubsec11.1.5\endcsname{11.1.5}
\expandafter\def\csname SGAsource@IV.subsubsec11.2.1\endcsname{11.2.1}
\expandafter\def\csname SGAsource@IV.subsubsec11.2.2\endcsname{11.2.2}
\expandafter\def\csname SGAsource@IV.subsubsec11.2.3\endcsname{11.2.3}
\expandafter\def\csname SGAsource@IV.subsubsec13.1.1\endcsname{13.1.1}
\expandafter\def\csname SGAsource@IV.subsubsec13.2.1\endcsname{13.2.1}
\expandafter\def\csname SGAsource@IV.subsubsec13.2.2\endcsname{13.2.2}
\expandafter\def\csname SGAsource@IV.subsubsec13.2.3\endcsname{13.2.3}
\expandafter\def\csname SGAsource@IV.subsubsec13.3.1\endcsname{13.3.1}
\expandafter\def\csname SGAsource@IV.subsubsec2.7.1\endcsname{2.7.1}
\expandafter\def\csname SGAsource@IV.subsubsec2.7.2\endcsname{2.7.2}
\expandafter\def\csname SGAsource@IV.subsubsec2.7.3\endcsname{2.7.3}
\expandafter\def\csname SGAsource@IV.subsubsec2.7.4\endcsname{2.7.4}
\expandafter\def\csname SGAsource@IV.subsubsec3.1.1\endcsname{3.1.1}
\expandafter\def\csname SGAsource@IV.subsubsec3.1.2\endcsname{3.1.2}
\expandafter\def\csname SGAsource@IV.subsubsec3.1.3\endcsname{3.1.3}
\expandafter\def\csname SGAsource@IV.subsubsec3.2.2\endcsname{3.2.2}
\expandafter\def\csname SGAsource@IV.subsubsec3.3.1\endcsname{3.3.1}
\expandafter\def\csname SGAsource@IV.subsubsec3.3.2\endcsname{3.3.2}
\expandafter\def\csname SGAsource@IV.subsubsec3.4.1\endcsname{3.4.1}
\expandafter\def\csname SGAsource@IV.subsubsec4.1.1\endcsname{4.1.1}
\expandafter\def\csname SGAsource@IV.subsubsec4.1.2\endcsname{4.1.2}
\expandafter\def\csname SGAsource@IV.subsubsec4.1.3\endcsname{4.1.3}
\expandafter\def\csname SGAsource@IV.subsubsec4.10.2\endcsname{4.10.2}
\expandafter\def\csname SGAsource@IV.subsubsec4.10.3\endcsname{4.10.3}
\expandafter\def\csname SGAsource@IV.subsubsec4.10.4\endcsname{4.10.4}
\expandafter\def\csname SGAsource@IV.subsubsec4.2.1\endcsname{4.2.1}
\expandafter\def\csname SGAsource@IV.subsubsec4.2.2\endcsname{4.2.2}
\expandafter\def\csname SGAsource@IV.subsubsec4.2.3\endcsname{4.2.3}
\expandafter\def\csname SGAsource@IV.subsubsec4.3.4\endcsname{4.3.4}
\expandafter\def\csname SGAsource@IV.subsubsec4.3.5\endcsname{4.3.5}
\expandafter\def\csname SGAsource@IV.subsubsec4.3.6\endcsname{4.3.6}
\expandafter\def\csname SGAsource@IV.subsubsec4.5.1\endcsname{4.5.1}
\expandafter\def\csname SGAsource@IV.subsubsec4.5.2\endcsname{4.5.2}
\expandafter\def\csname SGAsource@IV.subsubsec4.6.1\endcsname{4.6.1}
\expandafter\def\csname SGAsource@IV.subsubsec4.6.2\endcsname{4.6.2}
\expandafter\def\csname SGAsource@IV.subsubsec4.6.3\endcsname{4.6.3}
\expandafter\def\csname SGAsource@IV.subsubsec4.7.5\endcsname{4.7.5}
\expandafter\def\csname SGAsource@IV.subsubsec4.9.1\endcsname{4.9.1}
\expandafter\def\csname SGAsource@IV.subsubsec4.9.2\endcsname{4.9.2}
\expandafter\def\csname SGAsource@IV.subsubsec4.9.3\endcsname{4.9.3}
\expandafter\def\csname SGAsource@IV.subsubsec4.9.5\endcsname{4.9.5}
\expandafter\def\csname SGAsource@IV.subsubsec4.9.6\endcsname{4.9.6}
\expandafter\def\csname SGAsource@IV.subsubsec5.10.4\endcsname{5.10.4}
\expandafter\def\csname SGAsource@IV.subsubsec5.11.1\endcsname{5.11.1}
\expandafter\def\csname SGAsource@IV.subsubsec5.12.2\endcsname{5.12.2}
\expandafter\def\csname SGAsource@IV.subsubsec5.8.3\endcsname{5.8.3}
\expandafter\def\csname SGAsource@IV.subsubsec5.8.4\endcsname{5.8.4}
\expandafter\def\csname SGAsource@IV.subsubsec6.1.1\endcsname{6.1.1}
\expandafter\def\csname SGAsource@IV.subsubsec6.1.2\endcsname{6.1.2}
\expandafter\def\csname SGAsource@IV.subsubsec6.2.1\endcsname{6.2.1}
\expandafter\def\csname SGAsource@IV.subsubsec6.4.0\endcsname{6.4.0}
\expandafter\def\csname SGAsource@IV.subsubsec6.4.2\endcsname{6.4.2}
\expandafter\def\csname SGAsource@IV.subsubsec6.4.3\endcsname{6.4.3}
\expandafter\def\csname SGAsource@IV.subsubsec6.8.2\endcsname{6.8.2}
\expandafter\def\csname SGAsource@IV.subsubsec6.8.5\endcsname{6.8.5}
\expandafter\def\csname SGAsource@IV.subsubsec6.8.6\endcsname{6.8.6}
\expandafter\def\csname SGAsource@IV.subsubsec6.8.7\endcsname{6.8.7}
\expandafter\def\csname SGAsource@IV.subsubsec7.1.5\endcsname{7.1.5}
\expandafter\def\csname SGAsource@IV.subsubsec7.1.6\endcsname{7.1.6}
\expandafter\def\csname SGAsource@IV.subsubsec7.1.7\endcsname{7.1.7}
\expandafter\def\csname SGAsource@IV.subsubsec7.1.8\endcsname{7.1.8}
\expandafter\def\csname SGAsource@IV.subsubsec7.2.5\endcsname{7.2.5}
\expandafter\def\csname SGAsource@IV.subsubsec8.4.1\endcsname{8.4.1}
\expandafter\def\csname SGAsource@IV.subsubsec8.4.2\endcsname{8.4.2}
\expandafter\def\csname SGAsource@IV.subsubsec8.4.3\endcsname{8.4.3}
\expandafter\def\csname SGAsource@IV.subsubsec8.4.4\endcsname{8.4.4}
\expandafter\def\csname SGAsource@IV.subsubsec8.4.5\endcsname{8.4.5}
\expandafter\def\csname SGAsource@IV.subsubsec8.5.1\endcsname{8.5.1}
\expandafter\def\csname SGAsource@IV.subsubsec8.5.2\endcsname{8.5.2}
\expandafter\def\csname SGAsource@IV.subsubsec8.5.4\endcsname{8.5.4}
\expandafter\def\csname SGAsource@IV.subsubsec8.5.5\endcsname{8.5.5}
\expandafter\def\csname SGAsource@IV.subsubsec9.1.14\endcsname{9.1.14}
\expandafter\def\csname SGAsource@IV.subsubsec9.2.1\endcsname{9.2.1}
\expandafter\def\csname SGAsource@IV.subsubsec9.2.5\endcsname{9.2.5}
\expandafter\def\csname SGAsource@IV.subsubsec9.2.6\endcsname{9.2.6}
\expandafter\def\csname SGAsource@IV.subsubsec9.3.1\endcsname{9.3.1}
\expandafter\def\csname SGAsource@IV.subsubsec9.3.2\endcsname{9.3.2}
\expandafter\def\csname SGAsource@IV.subsubsec9.3.5\endcsname{9.3.5}
\expandafter\def\csname SGAsource@IV.subsubsec9.3.6\endcsname{9.3.6}
\expandafter\def\csname SGAsource@IV.subsubsec9.5.1\endcsname{9.5.1}
\expandafter\def\csname SGAsource@IV.subsubsec9.5.2\endcsname{9.5.2}
\expandafter\def\csname SGAsource@IV.subsubsec9.5.3\endcsname{9.5.3}
\expandafter\def\csname SGAsource@IV.subsubsec9.5.6\endcsname{9.5.6}
\expandafter\def\csname SGAsource@IV.theo1.2\endcsname{1.2}
\expandafter\def\csname SGAsource@IV.theo9.5.4\endcsname{9.5.4}
\expandafter\def\csname SGAsource@IX.Thr3.2\endcsname{3.2}
\expandafter\def\csname SGAsource@IX.Thr3.5\endcsname{3.5}
\expandafter\def\csname SGAsource@IX.cor2.7.1\endcsname{2.7.1}
\expandafter\def\csname SGAsource@IX.cor2.7.2\endcsname{2.7.2}
\expandafter\def\csname SGAsource@IX.cor2.7.3\endcsname{2.7.3}
\expandafter\def\csname SGAsource@IX.cor2.7.4\endcsname{2.7.4}
\expandafter\def\csname SGAsource@IX.cor3.4\endcsname{3.4}
\expandafter\def\csname SGAsource@IX.cor4.2\endcsname{4.2}
\expandafter\def\csname SGAsource@IX.cor4.7\endcsname{4.7}
\expandafter\def\csname SGAsource@IX.cor5.2\endcsname{5.2}
\expandafter\def\csname SGAsource@IX.cor5.7\endcsname{5.7}
\expandafter\def\csname SGAsource@IX.def1.1\endcsname{1.1}
\expandafter\def\csname SGAsource@IX.def1.3\endcsname{1.3}
\expandafter\def\csname SGAsource@IX.def1.5\endcsname{1.5}
\expandafter\def\csname SGAsource@IX.def2.12\endcsname{2.12}
\expandafter\def\csname SGAsource@IX.def2.3\endcsname{2.3}
\expandafter\def\csname SGAsource@IX.eq3.1\endcsname{3.1}
\expandafter\def\csname SGAsource@IX.eq4.3.1\endcsname{4.3.1}
\expandafter\def\csname SGAsource@IX.eq4.4.1\endcsname{4.4.1}
\expandafter\def\csname SGAsource@IX.eq4.4.2\endcsname{4.4.2}
\expandafter\def\csname SGAsource@IX.eq4.5\endcsname{4.5}
\expandafter\def\csname SGAsource@IX.eq4.8\endcsname{4.8}
\expandafter\def\csname SGAsource@IX.eq5.1.1\endcsname{5.1.1}
\expandafter\def\csname SGAsource@IX.eq5.1.2\endcsname{5.1.2}
\expandafter\def\csname SGAsource@IX.eq5.1.3\endcsname{5.1.3}
\expandafter\def\csname SGAsource@IX.eq5.1.4\endcsname{5.1.4}
\expandafter\def\csname SGAsource@IX.eq5.4\endcsname{5.4}
\expandafter\def\csname SGAsource@IX.lem2.1\endcsname{2.1}
\expandafter\def\csname SGAsource@IX.lem2.10\endcsname{2.10}
\expandafter\def\csname SGAsource@IX.lem2.14.1\endcsname{2.14.1}
\expandafter\def\csname SGAsource@IX.lem2.2\endcsname{2.2}
\expandafter\def\csname SGAsource@IX.lem2.8.1\endcsname{2.8.1}
\expandafter\def\csname SGAsource@IX.lem4.1\endcsname{4.1}
\expandafter\def\csname SGAsource@IX.par2.7.2.1\endcsname{2.7.2.1}
\expandafter\def\csname SGAsource@IX.prop1.2\endcsname{1.2}
\expandafter\def\csname SGAsource@IX.prop1.4\endcsname{1.4}
\expandafter\def\csname SGAsource@IX.prop1.6\endcsname{1.6}
\expandafter\def\csname SGAsource@IX.prop2.11\endcsname{2.11}
\expandafter\def\csname SGAsource@IX.prop2.13\endcsname{2.13}
\expandafter\def\csname SGAsource@IX.prop2.14\endcsname{2.14}
\expandafter\def\csname SGAsource@IX.prop2.4\endcsname{2.4}
\expandafter\def\csname SGAsource@IX.prop2.5\endcsname{2.5}
\expandafter\def\csname SGAsource@IX.prop2.6\endcsname{2.6}
\expandafter\def\csname SGAsource@IX.prop2.7\endcsname{2.7}
\expandafter\def\csname SGAsource@IX.prop2.8\endcsname{2.8}
\expandafter\def\csname SGAsource@IX.prop2.9\endcsname{2.9}
\expandafter\def\csname SGAsource@IX.prop3.6\endcsname{3.6}
\expandafter\def\csname SGAsource@IX.prop5.5\endcsname{5.5}
\expandafter\def\csname SGAsource@IX.prop5.6\endcsname{5.6}
\expandafter\def\csname SGAsource@IX.rem5.8\endcsname{5.8}
\expandafter\def\csname SGAsource@IX.rems2.14.2\endcsname{2.14.2}
\expandafter\def\csname SGAsource@IX.rems3.7\endcsname{3.7}
\expandafter\def\csname SGAsource@IX.sec0\endcsname{0}
\expandafter\def\csname SGAsource@IX.sec1\endcsname{1}
\expandafter\def\csname SGAsource@IX.sec2\endcsname{2}
\expandafter\def\csname SGAsource@IX.sec3\endcsname{3}
\expandafter\def\csname SGAsource@IX.sec4\endcsname{4}
\expandafter\def\csname SGAsource@IX.sec5\endcsname{5}
\expandafter\def\csname SGAsource@IX.subsec2.0\endcsname{2.0}
\expandafter\def\csname SGAsource@IX.subsec3.0\endcsname{3.0}
\expandafter\def\csname SGAsource@IX.subsec4.0\endcsname{4.0}
\expandafter\def\csname SGAsource@IX.subsec4.3\endcsname{4.3}
\expandafter\def\csname SGAsource@IX.subsec4.4\endcsname{4.4}
\expandafter\def\csname SGAsource@IX.subsec5.1\endcsname{5.1}
\expandafter\def\csname SGAsource@IX.subsubsec2.3.1\endcsname{2.3.1}
\expandafter\def\csname SGAsource@IX.subsubsec3.3.1\endcsname{3.3.1}
\expandafter\def\csname SGAsource@IX.subsubsec4.1.1\endcsname{4.1.1}
\expandafter\def\csname SGAsource@IX.subsubsec5.2.1\endcsname{5.2.1}
\expandafter\def\csname SGAsource@IX.thm3.3\endcsname{3.3}
\expandafter\def\csname SGAsource@IX.thm4.6\endcsname{4.6}
\expandafter\def\csname SGAsource@Ii.prop2.5.3\endcsname{2.5.3}
\expandafter\def\csname SGAsource@V.1.10.1\endcsname{1.10.1}
\expandafter\def\csname SGAsource@V.1.10.2\endcsname{1.10.2}
\expandafter\def\csname SGAsource@V.1.10.3\endcsname{1.10.3}
\expandafter\def\csname SGAsource@V.1.10.4\endcsname{1.10.4}
\expandafter\def\csname SGAsource@V.1.10.5\endcsname{1.10.5}
\expandafter\def\csname SGAsource@V.1.4.2\endcsname{1.4.2}
\expandafter\def\csname SGAsource@V.1.4.3\endcsname{1.4.3}
\expandafter\def\csname SGAsource@V.1.5.1\endcsname{1.5.1}
\expandafter\def\csname SGAsource@V.1.5.2\endcsname{1.5.2}
\expandafter\def\csname SGAsource@V.1.5.3\endcsname{1.5.3}
\expandafter\def\csname SGAsource@V.1.5.4\endcsname{1.5.4}
\expandafter\def\csname SGAsource@V.1.5.5\endcsname{1.5.5}
\expandafter\def\csname SGAsource@V.2.2.1\endcsname{2.2.1}
\expandafter\def\csname SGAsource@V.2.2.2\endcsname{2.2.2}
\expandafter\def\csname SGAsource@V.2.3.3.1\endcsname{2.3.3.1}
\expandafter\def\csname SGAsource@V.2.4.2.1\endcsname{2.4.2.1}
\expandafter\def\csname SGAsource@V.2.4.2.2\endcsname{2.4.2.2}
\expandafter\def\csname SGAsource@V.2.4.5.1\endcsname{2.4.5.1}
\expandafter\def\csname SGAsource@V.2.4.5.2\endcsname{2.4.5.2}
\expandafter\def\csname SGAsource@V.2.4.5.3\endcsname{2.4.5.3}
\expandafter\def\csname SGAsource@V.2.4.5.4\endcsname{2.4.5.4}
\expandafter\def\csname SGAsource@V.2.4.6.1\endcsname{2.4.6.1}
\expandafter\def\csname SGAsource@V.2.5.1.1\endcsname{2.5.1.1}
\expandafter\def\csname SGAsource@V.2.5.1.2\endcsname{2.5.1.2}
\expandafter\def\csname SGAsource@V.2.5.1.3\endcsname{2.5.1.3}
\expandafter\def\csname SGAsource@V.3.2.1\endcsname{3.2.1}
\expandafter\def\csname SGAsource@V.3.3.1\endcsname{3.3.1}
\expandafter\def\csname SGAsource@V.3.4.1\endcsname{3.4.1}
\expandafter\def\csname SGAsource@V.3.4.2\endcsname{3.4.2}
\expandafter\def\csname SGAsource@V.3.4.3\endcsname{3.4.3}
\expandafter\def\csname SGAsource@V.3.6.1\endcsname{3.6.1}
\expandafter\def\csname SGAsource@V.4.11.1\endcsname{4.11.1}
\expandafter\def\csname SGAsource@V.5.1.1\endcsname{5.1.1}
\expandafter\def\csname SGAsource@V.5.3.1\endcsname{5.3.1}
\expandafter\def\csname SGAsource@V.5.3.2\endcsname{5.3.2}
\expandafter\def\csname SGAsource@V.6.0.1\endcsname{6.0.1}
\expandafter\def\csname SGAsource@V.6.0.2\endcsname{6.0.2}
\expandafter\def\csname SGAsource@V.6.0.3\endcsname{6.0.3}
\expandafter\def\csname SGAsource@V.6.0.4\endcsname{6.0.4}
\expandafter\def\csname SGAsource@V.6.1.1\endcsname{6.1.1}
\expandafter\def\csname SGAsource@V.6.1.2\endcsname{6.1.2}
\expandafter\def\csname SGAsource@V.6.1.3\endcsname{6.1.3}
\expandafter\def\csname SGAsource@V.6.10.1\endcsname{6.10.1}
\expandafter\def\csname SGAsource@V.6.13.5\endcsname{6.13.5}
\expandafter\def\csname SGAsource@V.6.13.6\endcsname{6.13.6}
\expandafter\def\csname SGAsource@V.6.13.7\endcsname{6.13.7}
\expandafter\def\csname SGAsource@V.6.13.8\endcsname{6.13.8}
\expandafter\def\csname SGAsource@V.6.14.1\endcsname{6.14.1}
\expandafter\def\csname SGAsource@V.6.14.2\endcsname{6.14.2}
\expandafter\def\csname SGAsource@V.6.15.1\endcsname{6.15.1}
\expandafter\def\csname SGAsource@V.6.2.1\endcsname{6.2.1}
\expandafter\def\csname SGAsource@V.6.2.2\endcsname{6.2.2}
\expandafter\def\csname SGAsource@V.6.3.3\endcsname{6.3.3}
\expandafter\def\csname SGAsource@V.6.3.5\endcsname{6.3.5}
\expandafter\def\csname SGAsource@V.6.3.6\endcsname{6.3.6}
\expandafter\def\csname SGAsource@V.6.3.7\endcsname{6.3.7}
\expandafter\def\csname SGAsource@V.6.4.1\endcsname{6.4.1}
\expandafter\def\csname SGAsource@V.6.4.2\endcsname{6.4.2}
\expandafter\def\csname SGAsource@V.6.4.3\endcsname{6.4.3}
\expandafter\def\csname SGAsource@V.6.5.1\endcsname{6.5.1}
\expandafter\def\csname SGAsource@V.6.5.2\endcsname{6.5.2}
\expandafter\def\csname SGAsource@V.6.5.3\endcsname{6.5.3}
\expandafter\def\csname SGAsource@V.6.5.4\endcsname{6.5.4}
\expandafter\def\csname SGAsource@V.6.5.5\endcsname{6.5.5}
\expandafter\def\csname SGAsource@V.6.8.1\endcsname{6.8.1}
\expandafter\def\csname SGAsource@V.6.8.2\endcsname{6.8.2}
\expandafter\def\csname SGAsource@V.6.8.3\endcsname{6.8.3}
\expandafter\def\csname SGAsource@V.6.8.4\endcsname{6.8.4}
\expandafter\def\csname SGAsource@V.6.8.5\endcsname{6.8.5}
\expandafter\def\csname SGAsource@V.6.8.6\endcsname{6.8.6}
\expandafter\def\csname SGAsource@V.6.8.7\endcsname{6.8.7}
\expandafter\def\csname SGAsource@V.6.9.1\endcsname{6.9.1}
\expandafter\def\csname SGAsource@V.6.9.2\endcsname{6.9.2}
\expandafter\def\csname SGAsource@V.7.2.2\endcsname{7.2.2}
\expandafter\def\csname SGAsource@V.7.4.0.1\endcsname{7.4.0.1}
\expandafter\def\csname SGAsource@V.7.4.0.2\endcsname{7.4.0.2}
\expandafter\def\csname SGAsource@V.7.4.0.3\endcsname{7.4.0.3}
\expandafter\def\csname SGAsource@V.7.4.0.4\endcsname{7.4.0.4}
\expandafter\def\csname SGAsource@V.8.2.7.1\endcsname{8.2.7.1}
\expandafter\def\csname SGAsource@V.cor1.3.3\endcsname{1.3.3}
\expandafter\def\csname SGAsource@V.cor1.7.1\endcsname{1.7.1}
\expandafter\def\csname SGAsource@V.cor1.7.2\endcsname{1.7.2}
\expandafter\def\csname SGAsource@V.cor2.3.5\endcsname{2.3.5}
\expandafter\def\csname SGAsource@V.cor3.3\endcsname{3.3}
\expandafter\def\csname SGAsource@V.cor3.4\endcsname{3.4}
\expandafter\def\csname SGAsource@V.cor3.5\endcsname{3.5}
\expandafter\def\csname SGAsource@V.cor4.5\endcsname{4.5}
\expandafter\def\csname SGAsource@V.cor8.2.9\endcsname{8.2.9}
\expandafter\def\csname SGAsource@V.def1.1\endcsname{1.1}
\expandafter\def\csname SGAsource@V.def1.8\endcsname{1.8}
\expandafter\def\csname SGAsource@V.def4.1\endcsname{4.1}
\expandafter\def\csname SGAsource@V.def6.12\endcsname{6.12}
\expandafter\def\csname SGAsource@V.def7.1.1\endcsname{7.1.1}
\expandafter\def\csname SGAsource@V.def7.3.3\endcsname{7.3.3}
\expandafter\def\csname SGAsource@V.def8.1.1\endcsname{8.1.1}
\expandafter\def\csname SGAsource@V.def8.1.2.1\endcsname{8.1.2.1}
\expandafter\def\csname SGAsource@V.def8.1.3\endcsname{8.1.3}
\expandafter\def\csname SGAsource@V.def8.1.4\endcsname{8.1.4}
\expandafter\def\csname SGAsource@V.def8.2.1\endcsname{8.2.1}
\expandafter\def\csname SGAsource@V.eq1.4.1\endcsname{1.4.1}
\expandafter\def\csname SGAsource@V.exe0.5.2\endcsname{0.5.2}
\expandafter\def\csname SGAsource@V.exem6.12.3\endcsname{6.12.3}
\expandafter\def\csname SGAsource@V.exer2.3.6\endcsname{2.3.6}
\expandafter\def\csname SGAsource@V.exer3.6\endcsname{3.6}
\expandafter\def\csname SGAsource@V.exer4.13\endcsname{4.13}
\expandafter\def\csname SGAsource@V.exer4.15\endcsname{4.15}
\expandafter\def\csname SGAsource@V.exer4.16\endcsname{4.16}
\expandafter\def\csname SGAsource@V.exer4.17\endcsname{4.17}
\expandafter\def\csname SGAsource@V.lem1.6.1\endcsname{1.6.1}
\expandafter\def\csname SGAsource@V.lem7.2.1\endcsname{7.2.1}
\expandafter\def\csname SGAsource@V.lem7.3.4\endcsname{7.3.4}
\expandafter\def\csname SGAsource@V.lem7.3.6\endcsname{7.3.6}
\expandafter\def\csname SGAsource@V.lem8.1.11\endcsname{8.1.11}
\expandafter\def\csname SGAsource@V.lem8.1.6\endcsname{8.1.6}
\expandafter\def\csname SGAsource@V.lem8.1.8\endcsname{8.1.8}
\expandafter\def\csname SGAsource@V.lem8.2.10\endcsname{8.2.10}
\expandafter\def\csname SGAsource@V.lem8.2.11\endcsname{8.2.11}
\expandafter\def\csname SGAsource@V.lem8.2.3\endcsname{8.2.3}
\expandafter\def\csname SGAsource@V.lem8.2.8\endcsname{8.2.8}
\expandafter\def\csname SGAsource@V.par7.3.1.1\endcsname{7.3.1.1}
\expandafter\def\csname SGAsource@V.par7.3.1.2\endcsname{7.3.1.2}
\expandafter\def\csname SGAsource@V.par7.3.1.3\endcsname{7.3.1.3}
\expandafter\def\csname SGAsource@V.par7.3.1.4\endcsname{7.3.1.4}
\expandafter\def\csname SGAsource@V.par7.3.1.6\endcsname{7.3.1.6}
\expandafter\def\csname SGAsource@V.par7.3.1.7\endcsname{7.3.1.7}
\expandafter\def\csname SGAsource@V.par7.3.6.0\endcsname{7.3.6.0}
\expandafter\def\csname SGAsource@V.par8.1.1.1\endcsname{8.1.1.1}
\expandafter\def\csname SGAsource@V.par8.1.13.0\endcsname{8.1.13.0}
\expandafter\def\csname SGAsource@V.par8.1.9.0\endcsname{8.1.9.0}
\expandafter\def\csname SGAsource@V.prob4.14\endcsname{4.14}
\expandafter\def\csname SGAsource@V.prop0.1\endcsname{0.1}
\expandafter\def\csname SGAsource@V.prop0.2\endcsname{0.2}
\expandafter\def\csname SGAsource@V.prop0.3\endcsname{0.3}
\expandafter\def\csname SGAsource@V.prop0.4\endcsname{0.4}
\expandafter\def\csname SGAsource@V.prop0.5\endcsname{0.5}
\expandafter\def\csname SGAsource@V.prop1.11\endcsname{1.11}
\expandafter\def\csname SGAsource@V.prop1.2\endcsname{1.2}
\expandafter\def\csname SGAsource@V.prop1.3.1\endcsname{1.3.1}
\expandafter\def\csname SGAsource@V.prop1.3.2\endcsname{1.3.2}
\expandafter\def\csname SGAsource@V.prop1.5\endcsname{1.5}
\expandafter\def\csname SGAsource@V.prop1.6\endcsname{1.6}
\expandafter\def\csname SGAsource@V.prop1.7\endcsname{1.7}
\expandafter\def\csname SGAsource@V.prop1.9\endcsname{1.9}
\expandafter\def\csname SGAsource@V.prop2.3.4\endcsname{2.3.4}
\expandafter\def\csname SGAsource@V.prop3.1\endcsname{3.1}
\expandafter\def\csname SGAsource@V.prop4.10\endcsname{4.10}
\expandafter\def\csname SGAsource@V.prop4.11\endcsname{4.11}
\expandafter\def\csname SGAsource@V.prop4.3\endcsname{4.3}
\expandafter\def\csname SGAsource@V.prop4.7\endcsname{4.7}
\expandafter\def\csname SGAsource@V.prop4.9\endcsname{4.9}
\expandafter\def\csname SGAsource@V.prop5.1\endcsname{5.1}
\expandafter\def\csname SGAsource@V.prop5.2\endcsname{5.2}
\expandafter\def\csname SGAsource@V.prop5.3\endcsname{5.3}
\expandafter\def\csname SGAsource@V.prop5.4\endcsname{5.4}
\expandafter\def\csname SGAsource@V.prop6.1\endcsname{6.1}
\expandafter\def\csname SGAsource@V.prop6.10\endcsname{6.10}
\expandafter\def\csname SGAsource@V.prop6.15\endcsname{6.15}
\expandafter\def\csname SGAsource@V.prop6.2\endcsname{6.2}
\expandafter\def\csname SGAsource@V.prop6.4\endcsname{6.4}
\expandafter\def\csname SGAsource@V.prop6.5\endcsname{6.5}
\expandafter\def\csname SGAsource@V.prop6.6\endcsname{6.6}
\expandafter\def\csname SGAsource@V.prop6.7\endcsname{6.7}
\expandafter\def\csname SGAsource@V.prop6.9\endcsname{6.9}
\expandafter\def\csname SGAsource@V.prop7.1.2\endcsname{7.1.2}
\expandafter\def\csname SGAsource@V.prop8.1.12\endcsname{8.1.12}
\expandafter\def\csname SGAsource@V.prop8.1.13.1\endcsname{8.1.13.1}
\expandafter\def\csname SGAsource@V.prop8.1.14\endcsname{8.1.14}
\expandafter\def\csname SGAsource@V.prop8.1.5\endcsname{8.1.5}
\expandafter\def\csname SGAsource@V.prop8.1.7\endcsname{8.1.7}
\expandafter\def\csname SGAsource@V.prop8.1.9\endcsname{8.1.9}
\expandafter\def\csname SGAsource@V.prop8.2.4\endcsname{8.2.4}
\expandafter\def\csname SGAsource@V.prop8.2.5\endcsname{8.2.5}
\expandafter\def\csname SGAsource@V.prop8.2.6\endcsname{8.2.6}
\expandafter\def\csname SGAsource@V.prop8.2.7\endcsname{8.2.7}
\expandafter\def\csname SGAsource@V.rem0.2.1\endcsname{0.2.1}
\expandafter\def\csname SGAsource@V.rem0.5.1\endcsname{0.5.1}
\expandafter\def\csname SGAsource@V.rem7.4.2\endcsname{7.4.2}
\expandafter\def\csname SGAsource@V.sec0\endcsname{0}
\expandafter\def\csname SGAsource@V.sec1\endcsname{1}
\expandafter\def\csname SGAsource@V.sec2\endcsname{2}
\expandafter\def\csname SGAsource@V.sec3\endcsname{3}
\expandafter\def\csname SGAsource@V.sec4\endcsname{4}
\expandafter\def\csname SGAsource@V.sec5\endcsname{5}
\expandafter\def\csname SGAsource@V.sec6\endcsname{6}
\expandafter\def\csname SGAsource@V.sec7\endcsname{7}
\expandafter\def\csname SGAsource@V.sec8\endcsname{8}
\expandafter\def\csname SGAsource@V.subsec1.10\endcsname{1.10}
\expandafter\def\csname SGAsource@V.subsec1.4\endcsname{1.4}
\expandafter\def\csname SGAsource@V.subsec2.1\endcsname{2.1}
\expandafter\def\csname SGAsource@V.subsec2.2\endcsname{2.2}
\expandafter\def\csname SGAsource@V.subsec2.3\endcsname{2.3}
\expandafter\def\csname SGAsource@V.subsec2.4\endcsname{2.4}
\expandafter\def\csname SGAsource@V.subsec2.5\endcsname{2.5}
\expandafter\def\csname SGAsource@V.subsec4.12\endcsname{4.12}
\expandafter\def\csname SGAsource@V.subsec4.2\endcsname{4.2}
\expandafter\def\csname SGAsource@V.subsec4.4\endcsname{4.4}
\expandafter\def\csname SGAsource@V.subsec4.6\endcsname{4.6}
\expandafter\def\csname SGAsource@V.subsec4.8\endcsname{4.8}
\expandafter\def\csname SGAsource@V.subsec5.0\endcsname{5.0}
\expandafter\def\csname SGAsource@V.subsec6.0\endcsname{6.0}
\expandafter\def\csname SGAsource@V.subsec6.11\endcsname{6.11}
\expandafter\def\csname SGAsource@V.subsec6.13\endcsname{6.13}
\expandafter\def\csname SGAsource@V.subsec6.14\endcsname{6.14}
\expandafter\def\csname SGAsource@V.subsec6.16\endcsname{6.16}
\expandafter\def\csname SGAsource@V.subsec6.3\endcsname{6.3}
\expandafter\def\csname SGAsource@V.subsec6.8\endcsname{6.8}
\expandafter\def\csname SGAsource@V.subsec7.1\endcsname{7.1}
\expandafter\def\csname SGAsource@V.subsec7.2\endcsname{7.2}
\expandafter\def\csname SGAsource@V.subsec7.3\endcsname{7.3}
\expandafter\def\csname SGAsource@V.subsec7.4\endcsname{7.4}
\expandafter\def\csname SGAsource@V.subsec8.1\endcsname{8.1}
\expandafter\def\csname SGAsource@V.subsec8.2\endcsname{8.2}
\expandafter\def\csname SGAsource@V.subsubsec0.1.1\endcsname{0.1.1}
\expandafter\def\csname SGAsource@V.subsubsec0.5.2\endcsname{0.5.2}
\expandafter\def\csname SGAsource@V.subsubsec2.1.1\endcsname{2.1.1}
\expandafter\def\csname SGAsource@V.subsubsec2.1.2\endcsname{2.1.2}
\expandafter\def\csname SGAsource@V.subsubsec2.3.1\endcsname{2.3.1}
\expandafter\def\csname SGAsource@V.subsubsec2.3.2\endcsname{2.3.2}
\expandafter\def\csname SGAsource@V.subsubsec2.3.3\endcsname{2.3.3}
\expandafter\def\csname SGAsource@V.subsubsec2.4.1\endcsname{2.4.1}
\expandafter\def\csname SGAsource@V.subsubsec2.4.2\endcsname{2.4.2}
\expandafter\def\csname SGAsource@V.subsubsec2.4.3\endcsname{2.4.3}
\expandafter\def\csname SGAsource@V.subsubsec2.4.4\endcsname{2.4.4}
\expandafter\def\csname SGAsource@V.subsubsec2.4.5\endcsname{2.4.5}
\expandafter\def\csname SGAsource@V.subsubsec2.4.6\endcsname{2.4.6}
\expandafter\def\csname SGAsource@V.subsubsec2.5.1\endcsname{2.5.1}
\expandafter\def\csname SGAsource@V.subsubsec2.5.2\endcsname{2.5.2}
\expandafter\def\csname SGAsource@V.subsubsec6.12.1\endcsname{6.12.1}
\expandafter\def\csname SGAsource@V.subsubsec6.12.2\endcsname{6.12.2}
\expandafter\def\csname SGAsource@V.subsubsec7.1.0\endcsname{7.1.0}
\expandafter\def\csname SGAsource@V.subsubsec7.2.0\endcsname{7.2.0}
\expandafter\def\csname SGAsource@V.subsubsec7.3.0\endcsname{7.3.0}
\expandafter\def\csname SGAsource@V.subsubsec7.3.1\endcsname{7.3.1}
\expandafter\def\csname SGAsource@V.subsubsec7.3.5\endcsname{7.3.5}
\expandafter\def\csname SGAsource@V.subsubsec7.3.7\endcsname{7.3.7}
\expandafter\def\csname SGAsource@V.subsubsec7.4.0\endcsname{7.4.0}
\expandafter\def\csname SGAsource@V.subsubsec8.1.0\endcsname{8.1.0}
\expandafter\def\csname SGAsource@V.subsubsec8.1.10\endcsname{8.1.10}
\expandafter\def\csname SGAsource@V.subsubsec8.1.2\endcsname{8.1.2}
\expandafter\def\csname SGAsource@V.subsubsec8.2.0\endcsname{8.2.0}
\expandafter\def\csname SGAsource@V.subusbsec1.6.2\endcsname{1.6.2}
\expandafter\def\csname SGAsource@V.thm3.2\endcsname{3.2}
\expandafter\def\csname SGAsource@V.thm7.3.2\endcsname{7.3.2}
\expandafter\def\csname SGAsource@V.thm7.4.1\endcsname{7.4.1}
\expandafter\def\csname SGAsource@V.thm8.2.12\endcsname{8.2.12}
\expandafter\def\csname SGAsource@V.thm8.2.2\endcsname{8.2.2}
\expandafter\def\csname SGAsource@VI.cor1.11\endcsname{1.11}
\expandafter\def\csname SGAsource@VI.cor1.12\endcsname{1.12}
\expandafter\def\csname SGAsource@VI.cor1.15\endcsname{1.15}
\expandafter\def\csname SGAsource@VI.cor1.16\endcsname{1.16}
\expandafter\def\csname SGAsource@VI.cor1.17\endcsname{1.17}
\expandafter\def\csname SGAsource@VI.cor1.17.1\endcsname{1.17.1}
\expandafter\def\csname SGAsource@VI.cor1.18\endcsname{1.18}
\expandafter\def\csname SGAsource@VI.cor1.19\endcsname{1.19}
\expandafter\def\csname SGAsource@VI.cor1.24\endcsname{1.24}
\expandafter\def\csname SGAsource@VI.cor1.24.2\endcsname{1.24.2}
\expandafter\def\csname SGAsource@VI.cor1.25\endcsname{1.25}
\expandafter\def\csname SGAsource@VI.cor1.26\endcsname{1.26}
\expandafter\def\csname SGAsource@VI.cor1.4\endcsname{1.4}
\expandafter\def\csname SGAsource@VI.cor1.8.1\endcsname{1.8.1}
\expandafter\def\csname SGAsource@VI.cor1.8.2\endcsname{1.8.2}
\expandafter\def\csname SGAsource@VI.cor2.1.1\endcsname{2.1.1}
\expandafter\def\csname SGAsource@VI.cor2.1.2\endcsname{2.1.2}
\expandafter\def\csname SGAsource@VI.cor2.6\endcsname{2.6}
\expandafter\def\csname SGAsource@VI.cor2.7\endcsname{2.7}
\expandafter\def\csname SGAsource@VI.cor2.8\endcsname{2.8}
\expandafter\def\csname SGAsource@VI.cor2.9.4\endcsname{2.9.4}
\expandafter\def\csname SGAsource@VI.cor3.3\endcsname{3.3}
\expandafter\def\csname SGAsource@VI.cor3.6\endcsname{3.6}
\expandafter\def\csname SGAsource@VI.cor3.7\endcsname{3.7}
\expandafter\def\csname SGAsource@VI.cor3.8\endcsname{3.8}
\expandafter\def\csname SGAsource@VI.cor4.11\endcsname{4.11}
\expandafter\def\csname SGAsource@VI.cor4.14\endcsname{4.14}
\expandafter\def\csname SGAsource@VI.cor4.3\endcsname{4.3}
\expandafter\def\csname SGAsource@VI.cor4.4\endcsname{4.4}
\expandafter\def\csname SGAsource@VI.cor4.5\endcsname{4.5}
\expandafter\def\csname SGAsource@VI.cor4.7\endcsname{4.7}
\expandafter\def\csname SGAsource@VI.cor5.10\endcsname{5.10}
\expandafter\def\csname SGAsource@VI.cor5.2\endcsname{5.2}
\expandafter\def\csname SGAsource@VI.cor5.3\endcsname{5.3}
\expandafter\def\csname SGAsource@VI.cor5.5\endcsname{5.5}
\expandafter\def\csname SGAsource@VI.cor5.8\endcsname{5.8}
\expandafter\def\csname SGAsource@VI.cor7.4.11\endcsname{7.4.11}
\expandafter\def\csname SGAsource@VI.cor8.3.14\endcsname{8.3.14}
\expandafter\def\csname SGAsource@VI.cor8.5.10\endcsname{8.5.10}
\expandafter\def\csname SGAsource@VI.cor8.5.4\endcsname{8.5.4}
\expandafter\def\csname SGAsource@VI.cor8.5.5\endcsname{8.5.5}
\expandafter\def\csname SGAsource@VI.cor8.5.6\endcsname{8.5.6}
\expandafter\def\csname SGAsource@VI.cor8.5.7\endcsname{8.5.7}
\expandafter\def\csname SGAsource@VI.cor8.7.4\endcsname{8.7.4}
\expandafter\def\csname SGAsource@VI.cor8.7.5\endcsname{8.7.5}
\expandafter\def\csname SGAsource@VI.cor8.7.6\endcsname{8.7.6}
\expandafter\def\csname SGAsource@VI.cor8.7.7\endcsname{8.7.7}
\expandafter\def\csname SGAsource@VI.cor8.7.9\endcsname{8.7.9}
\expandafter\def\csname SGAsource@VI.cor9.5\endcsname{9.5}
\expandafter\def\csname SGAsource@VI.def1.1\endcsname{1.1}
\expandafter\def\csname SGAsource@VI.def1.13\endcsname{1.13}
\expandafter\def\csname SGAsource@VI.def1.30\endcsname{1.30}
\expandafter\def\csname SGAsource@VI.def1.7\endcsname{1.7}
\expandafter\def\csname SGAsource@VI.def2.11\endcsname{2.11}
\expandafter\def\csname SGAsource@VI.def2.3\endcsname{2.3}
\expandafter\def\csname SGAsource@VI.def2.9.1\endcsname{2.9.1}
\expandafter\def\csname SGAsource@VI.def3.1\endcsname{3.1}
\expandafter\def\csname SGAsource@VI.def3.7\endcsname{3.7}
\expandafter\def\csname SGAsource@VI.def5.6\endcsname{5.6}
\expandafter\def\csname SGAsource@VI.def6.1\endcsname{6.1}
\expandafter\def\csname SGAsource@VI.def6.3\endcsname{6.3}
\expandafter\def\csname SGAsource@VI.def7.1.1\endcsname{7.1.1}
\expandafter\def\csname SGAsource@VI.def7.1.5\endcsname{7.1.5}
\expandafter\def\csname SGAsource@VI.def7.4.1\endcsname{7.4.1}
\expandafter\def\csname SGAsource@VI.def8.1.1\endcsname{8.1.1}
\expandafter\def\csname SGAsource@VI.def8.2.5\endcsname{8.2.5}
\expandafter\def\csname SGAsource@VI.eq1.23.1\endcsname{1.23.1}
\expandafter\def\csname SGAsource@VI.eq1.24.1\endcsname{1.24.1}
\expandafter\def\csname SGAsource@VI.eq4.10.1\endcsname{4.10.1}
\expandafter\def\csname SGAsource@VI.eq4.10.2\endcsname{4.10.2}
\expandafter\def\csname SGAsource@VI.eq4.10.3\endcsname{4.10.3}
\expandafter\def\csname SGAsource@VI.eq4.10.4\endcsname{4.10.4}
\expandafter\def\csname SGAsource@VI.eq4.10.5\endcsname{4.10.5}
\expandafter\def\csname SGAsource@VI.eq4.10.6\endcsname{4.10.6}
\expandafter\def\csname SGAsource@VI.eq4.13.1\endcsname{4.13.1}
\expandafter\def\csname SGAsource@VI.eq4.13.2\endcsname{4.13.2}
\expandafter\def\csname SGAsource@VI.eq4.13.3\endcsname{4.13.3}
\expandafter\def\csname SGAsource@VI.eq4.13.4\endcsname{4.13.4}
\expandafter\def\csname SGAsource@VI.eq4.9.1\endcsname{4.9.1}
\expandafter\def\csname SGAsource@VI.eq4.9.2\endcsname{4.9.2}
\expandafter\def\csname SGAsource@VI.eq4.9.4\endcsname{4.9.4}
\expandafter\def\csname SGAsource@VI.eq4.9.5\endcsname{4.9.5}
\expandafter\def\csname SGAsource@VI.eq4.9.6\endcsname{4.9.6}
\expandafter\def\csname SGAsource@VI.eq4.9.7\endcsname{4.9.7}
\expandafter\def\csname SGAsource@VI.eq5.5.1\endcsname{5.5.1}
\expandafter\def\csname SGAsource@VI.eq5.9.1\endcsname{5.9.1}
\expandafter\def\csname SGAsource@VI.eq5.9.2\endcsname{5.9.2}
\expandafter\def\csname SGAsource@VI.eq6.1.2.1\endcsname{6.1.2.1}
\expandafter\def\csname SGAsource@VI.eq6.1.2.2\endcsname{6.1.2.2}
\expandafter\def\csname SGAsource@VI.eq7.1.3.1\endcsname{7.1.3.1}
\expandafter\def\csname SGAsource@VI.eq7.1.3.2\endcsname{7.1.3.2}
\expandafter\def\csname SGAsource@VI.eq7.1.6.1\endcsname{7.1.6.1}
\expandafter\def\csname SGAsource@VI.eq7.1.6.2\endcsname{7.1.6.2}
\expandafter\def\csname SGAsource@VI.eq7.1.6.3\endcsname{7.1.6.3}
\expandafter\def\csname SGAsource@VI.eq7.1.6.4\endcsname{7.1.6.4}
\expandafter\def\csname SGAsource@VI.eq7.1.9.1\endcsname{7.1.9.1}
\expandafter\def\csname SGAsource@VI.eq7.1.9.2\endcsname{7.1.9.2}
\expandafter\def\csname SGAsource@VI.eq7.2.2.1\endcsname{7.2.2.1}
\expandafter\def\csname SGAsource@VI.eq7.2.2.2\endcsname{7.2.2.2}
\expandafter\def\csname SGAsource@VI.eq7.2.5.1\endcsname{7.2.5.1}
\expandafter\def\csname SGAsource@VI.eq7.2.6.1\endcsname{7.2.6.1}
\expandafter\def\csname SGAsource@VI.eq7.2.6.2\endcsname{7.2.6.2}
\expandafter\def\csname SGAsource@VI.eq7.2.6.3\endcsname{7.2.6.3}
\expandafter\def\csname SGAsource@VI.eq7.2.6.4\endcsname{7.2.6.4}
\expandafter\def\csname SGAsource@VI.eq7.2.6.5\endcsname{7.2.6.5}
\expandafter\def\csname SGAsource@VI.eq7.2.6.6\endcsname{7.2.6.6}
\expandafter\def\csname SGAsource@VI.eq7.2.6.7\endcsname{7.2.6.7}
\expandafter\def\csname SGAsource@VI.eq7.3.1.1\endcsname{7.3.1.1}
\expandafter\def\csname SGAsource@VI.eq7.3.1.2\endcsname{7.3.1.2}
\expandafter\def\csname SGAsource@VI.eq7.3.1.3\endcsname{7.3.1.3}
\expandafter\def\csname SGAsource@VI.eq7.3.1.4\endcsname{7.3.1.4}
\expandafter\def\csname SGAsource@VI.eq7.3.1.5\endcsname{7.3.1.5}
\expandafter\def\csname SGAsource@VI.eq7.3.2.1\endcsname{7.3.2.1}
\expandafter\def\csname SGAsource@VI.eq7.3.2.2\endcsname{7.3.2.2}
\expandafter\def\csname SGAsource@VI.eq7.3.2.3\endcsname{7.3.2.3}
\expandafter\def\csname SGAsource@VI.eq7.3.2.4\endcsname{7.3.2.4}
\expandafter\def\csname SGAsource@VI.eq7.3.2.5\endcsname{7.3.2.5}
\expandafter\def\csname SGAsource@VI.eq7.3.2.6\endcsname{7.3.2.6}
\expandafter\def\csname SGAsource@VI.eq7.3.2.7\endcsname{7.3.2.7}
\expandafter\def\csname SGAsource@VI.eq7.4.10.1\endcsname{7.4.10.1}
\expandafter\def\csname SGAsource@VI.eq7.4.10.2\endcsname{7.4.10.2}
\expandafter\def\csname SGAsource@VI.eq7.4.10.3\endcsname{7.4.10.3}
\expandafter\def\csname SGAsource@VI.eq7.4.10.4\endcsname{7.4.10.4}
\expandafter\def\csname SGAsource@VI.eq7.4.10.5\endcsname{7.4.10.5}
\expandafter\def\csname SGAsource@VI.eq7.4.10.6\endcsname{7.4.10.6}
\expandafter\def\csname SGAsource@VI.eq7.4.5.1\endcsname{7.4.5.1}
\expandafter\def\csname SGAsource@VI.eq7.4.5.2\endcsname{7.4.5.2}
\expandafter\def\csname SGAsource@VI.eq7.4.5.3\endcsname{7.4.5.3}
\expandafter\def\csname SGAsource@VI.eq7.4.5.4\endcsname{7.4.5.4}
\expandafter\def\csname SGAsource@VI.eq7.4.5.5\endcsname{7.4.5.5}
\expandafter\def\csname SGAsource@VI.eq7.4.5.6\endcsname{7.4.5.6}
\expandafter\def\csname SGAsource@VI.eq7.4.5.7\endcsname{7.4.5.7}
\expandafter\def\csname SGAsource@VI.eq7.4.5.8\endcsname{7.4.5.8}
\expandafter\def\csname SGAsource@VI.eq7.4.6.1\endcsname{7.4.6.1}
\expandafter\def\csname SGAsource@VI.eq7.4.9.1\endcsname{7.4.9.1}
\expandafter\def\csname SGAsource@VI.eq7.4.9.2\endcsname{7.4.9.2}
\expandafter\def\csname SGAsource@VI.eq8.1.1.1\endcsname{8.1.1.1}
\expandafter\def\csname SGAsource@VI.eq8.1.2.1\endcsname{8.1.2.1}
\expandafter\def\csname SGAsource@VI.eq8.1.2.2\endcsname{8.1.2.2}
\expandafter\def\csname SGAsource@VI.eq8.1.3.1\endcsname{8.1.3.1.}
\expandafter\def\csname SGAsource@VI.eq8.1.4.1\endcsname{8.1.4.1}
\expandafter\def\csname SGAsource@VI.eq8.1.4.2\endcsname{8.1.4.2}
\expandafter\def\csname SGAsource@VI.eq8.1.6.1\endcsname{8.1.6.1}
\expandafter\def\csname SGAsource@VI.eq8.1.6.2\endcsname{8.1.6.2}
\expandafter\def\csname SGAsource@VI.eq8.1.6.3\endcsname{8.1.6.3}
\expandafter\def\csname SGAsource@VI.eq8.1.6.4\endcsname{8.1.6.4}
\expandafter\def\csname SGAsource@VI.eq8.1.6.5\endcsname{8.1.6.5}
\expandafter\def\csname SGAsource@VI.eq8.2.11.1\endcsname{8.2.11.1}
\expandafter\def\csname SGAsource@VI.eq8.2.8.1\endcsname{8.2.8.1}
\expandafter\def\csname SGAsource@VI.eq8.3.11.1\endcsname{8.3.11.1}
\expandafter\def\csname SGAsource@VI.eq8.3.13.1\endcsname{8.3.13.1}
\expandafter\def\csname SGAsource@VI.eq8.4.1.1\endcsname{8.4.1.1}
\expandafter\def\csname SGAsource@VI.eq8.4.2.1\endcsname{8.4.2.1}
\expandafter\def\csname SGAsource@VI.eq8.4.3.1\endcsname{8.4.3.1}
\expandafter\def\csname SGAsource@VI.eq8.5.1.1\endcsname{8.5.1.1}
\expandafter\def\csname SGAsource@VI.eq8.5.1.2\endcsname{8.5.1.2}
\expandafter\def\csname SGAsource@VI.eq8.5.1.3\endcsname{8.5.1.3}
\expandafter\def\csname SGAsource@VI.eq8.5.1.4\endcsname{8.5.1.4}
\expandafter\def\csname SGAsource@VI.eq8.5.1.5\endcsname{8.5.1.5}
\expandafter\def\csname SGAsource@VI.eq8.5.1.6\endcsname{8.5.1.6}
\expandafter\def\csname SGAsource@VI.eq8.5.1.7\endcsname{8.5.1.7}
\expandafter\def\csname SGAsource@VI.eq8.5.10.1\endcsname{8.5.10.1}
\expandafter\def\csname SGAsource@VI.eq8.5.2.1\endcsname{8.5.2.1}
\expandafter\def\csname SGAsource@VI.eq8.5.2.2\endcsname{8.5.2.2}
\expandafter\def\csname SGAsource@VI.eq8.5.2.3\endcsname{8.5.2.3}
\expandafter\def\csname SGAsource@VI.eq8.5.2.4\endcsname{8.5.2.4}
\expandafter\def\csname SGAsource@VI.eq8.5.2.5\endcsname{8.5.2.5}
\expandafter\def\csname SGAsource@VI.eq8.5.2.6\endcsname{8.5.2.6}
\expandafter\def\csname SGAsource@VI.eq8.5.2.7\endcsname{8.5.2.7}
\expandafter\def\csname SGAsource@VI.eq8.5.2.8\endcsname{8.5.2.8}
\expandafter\def\csname SGAsource@VI.eq8.5.2.9\endcsname{8.5.2.9}
\expandafter\def\csname SGAsource@VI.eq8.5.3.1\endcsname{8.5.3.1}
\expandafter\def\csname SGAsource@VI.eq8.5.5.1\endcsname{8.5.5.1}
\expandafter\def\csname SGAsource@VI.eq8.5.7.1\endcsname{8.5.7.1}
\expandafter\def\csname SGAsource@VI.eq8.5.7.2\endcsname{8.5.7.2}
\expandafter\def\csname SGAsource@VI.eq8.5.8.1\endcsname{8.5.8.1}
\expandafter\def\csname SGAsource@VI.eq8.5.8.2\endcsname{8.5.8.2}
\expandafter\def\csname SGAsource@VI.eq8.5.9.1\endcsname{8.5.9.1}
\expandafter\def\csname SGAsource@VI.eq8.7.3.1\endcsname{8.7.3.1}
\expandafter\def\csname SGAsource@VI.eq8.7.4.1\endcsname{8.7.4.1}
\expandafter\def\csname SGAsource@VI.eq8.7.5.1\endcsname{8.7.5.1}
\expandafter\def\csname SGAsource@VI.eq8.7.6.1\endcsname{8.7.6.1}
\expandafter\def\csname SGAsource@VI.eq8.7.6.2\endcsname{8.7.6.2}
\expandafter\def\csname SGAsource@VI.eq8.7.6.3\endcsname{8.7.6.3}
\expandafter\def\csname SGAsource@VI.eq8.7.7.1\endcsname{8.7.7.1}
\expandafter\def\csname SGAsource@VI.eq8.7.7.2\endcsname{8.7.7.2}
\expandafter\def\csname SGAsource@VI.eq8.7.9.1\endcsname{8.7.9.1}
\expandafter\def\csname SGAsource@VI.exe6.6\endcsname{6.6}
\expandafter\def\csname SGAsource@VI.exe6.8\endcsname{6.8}
\expandafter\def\csname SGAsource@VI.exe7.3.3\endcsname{7.3.3}
\expandafter\def\csname SGAsource@VI.exe7.4.14\endcsname{7.4.14}
\expandafter\def\csname SGAsource@VI.exe7.4.15\endcsname{7.4.15}
\expandafter\def\csname SGAsource@VI.exe8.3.15\endcsname{8.3.15}
\expandafter\def\csname SGAsource@VI.exem1.33\endcsname{1.33}
\expandafter\def\csname SGAsource@VI.exem1.6.1\endcsname{1.6.1}
\expandafter\def\csname SGAsource@VI.exem2.15\endcsname{2.15}
\expandafter\def\csname SGAsource@VI.exem3.10\endcsname{3.10}
\expandafter\def\csname SGAsource@VI.exem7.3.1\endcsname{7.3.1}
\expandafter\def\csname SGAsource@VI.exem7.3.2\endcsname{7.3.2}
\expandafter\def\csname SGAsource@VI.exer1.28\endcsname{1.28}
\expandafter\def\csname SGAsource@VI.exer2.16\endcsname{2.16}
\expandafter\def\csname SGAsource@VI.exer2.17\endcsname{2.17}
\expandafter\def\csname SGAsource@VI.exer2.18\endcsname{2.18}
\expandafter\def\csname SGAsource@VI.exer3.11\endcsname{3.11}
\expandafter\def\csname SGAsource@VI.exer3.12\endcsname{3.12}
\expandafter\def\csname SGAsource@VI.exer4.8\endcsname{4.8}
\expandafter\def\csname SGAsource@VI.lem1.25.1\endcsname{1.25.1}
\expandafter\def\csname SGAsource@VI.lem2.1.3\endcsname{2.1.3}
\expandafter\def\csname SGAsource@VI.lem4.12.1\endcsname{4.12.1}
\expandafter\def\csname SGAsource@VI.lem4.2.1\endcsname{4.2.1}
\expandafter\def\csname SGAsource@VI.lem4.5.1\endcsname{4.5.1}
\expandafter\def\csname SGAsource@VI.lem5.4\endcsname{5.4}
\expandafter\def\csname SGAsource@VI.lem7.4.12\endcsname{7.4.12}
\expandafter\def\csname SGAsource@VI.lem7.4.2.2\endcsname{7.4.2.2}
\expandafter\def\csname SGAsource@VI.lem8.2.1\endcsname{8.2.1}
\expandafter\def\csname SGAsource@VI.lem8.2.2\endcsname{8.2.2}
\expandafter\def\csname SGAsource@VI.lem8.3.11\endcsname{8.3.11}
\expandafter\def\csname SGAsource@VI.lem8.3.4\endcsname{8.3.4}
\expandafter\def\csname SGAsource@VI.lem8.7.2\endcsname{8.7.2}
\expandafter\def\csname SGAsource@VI.lem9.1\endcsname{9.1}
\expandafter\def\csname SGAsource@VI.lem9.2\endcsname{9.2}
\expandafter\def\csname SGAsource@VI.lem9.3\endcsname{9.3}
\expandafter\def\csname SGAsource@VI.lem9.4\endcsname{9.4}
\expandafter\def\csname SGAsource@VI.par7.2.6.8\endcsname{7.2.6.8}
\expandafter\def\csname SGAsource@VI.par7.4.13.2\endcsname{7.4.13.2}
\expandafter\def\csname SGAsource@VI.par7.4.13.3\endcsname{7.4.13.3}
\expandafter\def\csname SGAsource@VI.par7.4.13.4\endcsname{7.4.13.4}
\expandafter\def\csname SGAsource@VI.par7.4.13.5\endcsname{7.4.13.5}
\expandafter\def\csname SGAsource@VI.par7.4.13.6\endcsname{7.4.13.6}
\expandafter\def\csname SGAsource@VI.par7.4.13.7\endcsname{7.4.13.7}
\expandafter\def\csname SGAsource@VI.par7.4.2.1\endcsname{7.4.2.1}
\expandafter\def\csname SGAsource@VI.par8.1.3.2\endcsname{8.1.3.2}
\expandafter\def\csname SGAsource@VI.par8.1.3.3\endcsname{8.1.3.3}
\expandafter\def\csname SGAsource@VI.prob2.9.5\endcsname{2.9.5}
\expandafter\def\csname SGAsource@VI.prop1.10\endcsname{1.10}
\expandafter\def\csname SGAsource@VI.prop1.14\endcsname{1.14}
\expandafter\def\csname SGAsource@VI.prop1.2\endcsname{1.2}
\expandafter\def\csname SGAsource@VI.prop1.3\endcsname{1.3}
\expandafter\def\csname SGAsource@VI.prop1.8\endcsname{1.8}
\expandafter\def\csname SGAsource@VI.prop2.1\endcsname{2.1}
\expandafter\def\csname SGAsource@VI.prop2.10\endcsname{2.10}
\expandafter\def\csname SGAsource@VI.prop2.2\endcsname{2.2}
\expandafter\def\csname SGAsource@VI.prop2.5\endcsname{2.5}
\expandafter\def\csname SGAsource@VI.prop2.9\endcsname{2.9}
\expandafter\def\csname SGAsource@VI.prop2.9.3\endcsname{2.9.3}
\expandafter\def\csname SGAsource@VI.prop3.2\endcsname{3.2}
\expandafter\def\csname SGAsource@VI.prop3.4\endcsname{3.4}
\expandafter\def\csname SGAsource@VI.prop3.5\endcsname{3.5}
\expandafter\def\csname SGAsource@VI.prop4.10\endcsname{4.10}
\expandafter\def\csname SGAsource@VI.prop4.2\endcsname{4.2}
\expandafter\def\csname SGAsource@VI.prop4.6\endcsname{4.6}
\expandafter\def\csname SGAsource@VI.prop5.7\endcsname{5.7}
\expandafter\def\csname SGAsource@VI.prop6.2\endcsname{6.2}
\expandafter\def\csname SGAsource@VI.prop6.4\endcsname{6.4}
\expandafter\def\csname SGAsource@VI.prop6.5\endcsname{6.5}
\expandafter\def\csname SGAsource@VI.prop6.7\endcsname{6.7}
\expandafter\def\csname SGAsource@VI.prop6.9\endcsname{6.9}
\expandafter\def\csname SGAsource@VI.prop7.2.7\endcsname{7.2.7}
\expandafter\def\csname SGAsource@VI.prop7.4.10\endcsname{7.4.10}
\expandafter\def\csname SGAsource@VI.prop7.4.2\endcsname{7.4.2}
\expandafter\def\csname SGAsource@VI.prop7.4.4\endcsname{7.4.4}
\expandafter\def\csname SGAsource@VI.prop7.4.7\endcsname{7.4.7}
\expandafter\def\csname SGAsource@VI.prop8.1.7\endcsname{8.1.7}
\expandafter\def\csname SGAsource@VI.prop8.3.3\endcsname{8.3.3}
\expandafter\def\csname SGAsource@VI.prop8.3.6\endcsname{8.3.6}
\expandafter\def\csname SGAsource@VI.prop8.5.2\endcsname{8.5.2}
\expandafter\def\csname SGAsource@VI.prop8.5.3\endcsname{8.5.3}
\expandafter\def\csname SGAsource@VI.prop8.5.8\endcsname{8.5.8}
\expandafter\def\csname SGAsource@VI.prop8.5.9\endcsname{8.5.9}
\expandafter\def\csname SGAsource@VI.prop9.0\endcsname{9.0}
\expandafter\def\csname SGAsource@VI.rem1.20.1\endcsname{1.20.1}
\expandafter\def\csname SGAsource@VI.rem1.23.2\endcsname{1.23.2}
\expandafter\def\csname SGAsource@VI.rem1.27\endcsname{1.27}
\expandafter\def\csname SGAsource@VI.rem1.5.1\endcsname{1.5.1}
\expandafter\def\csname SGAsource@VI.rem1.9.1\endcsname{1.9.1}
\expandafter\def\csname SGAsource@VI.rem2.2.1\endcsname{2.2.1}
\expandafter\def\csname SGAsource@VI.rem2.8.1\endcsname{2.8.1}
\expandafter\def\csname SGAsource@VI.rem3.9\endcsname{3.9}
\expandafter\def\csname SGAsource@VI.rem4.13\endcsname{4.13}
\expandafter\def\csname SGAsource@VI.rem4.15\endcsname{4.15}
\expandafter\def\csname SGAsource@VI.rem4.6.1\endcsname{4.6.1}
\expandafter\def\csname SGAsource@VI.rem7.4.11.1\endcsname{7.4.11.1}
\expandafter\def\csname SGAsource@VI.rem7.4.3\endcsname{7.4.3}
\expandafter\def\csname SGAsource@VI.rem7.4.8\endcsname{7.4.8}
\expandafter\def\csname SGAsource@VI.rem8.2.4\endcsname{8.2.4}
\expandafter\def\csname SGAsource@VI.rem8.5.11\endcsname{8.5.11}
\expandafter\def\csname SGAsource@VI.rem8.7.8\endcsname{8.7.8}
\expandafter\def\csname SGAsource@VI.sec0\endcsname{0}
\expandafter\def\csname SGAsource@VI.sec1\endcsname{1}
\expandafter\def\csname SGAsource@VI.sec2\endcsname{2}
\expandafter\def\csname SGAsource@VI.sec3\endcsname{3}
\expandafter\def\csname SGAsource@VI.sec4\endcsname{4}
\expandafter\def\csname SGAsource@VI.sec5\endcsname{5}
\expandafter\def\csname SGAsource@VI.sec6\endcsname{6}
\expandafter\def\csname SGAsource@VI.sec7\endcsname{7}
\expandafter\def\csname SGAsource@VI.sec8\endcsname{8}
\expandafter\def\csname SGAsource@VI.sec9\endcsname{9}
\expandafter\def\csname SGAsource@VI.subsec1.21\endcsname{1.21}
\expandafter\def\csname SGAsource@VI.subsec1.22\endcsname{1.22}
\expandafter\def\csname SGAsource@VI.subsec1.31\endcsname{1.31}
\expandafter\def\csname SGAsource@VI.subsec1.32\endcsname{1.32}
\expandafter\def\csname SGAsource@VI.subsec2.12\endcsname{2.12}
\expandafter\def\csname SGAsource@VI.subsec2.13\endcsname{2.13}
\expandafter\def\csname SGAsource@VI.subsec2.14\endcsname{2.14}
\expandafter\def\csname SGAsource@VI.subsec2.4\endcsname{2.4}
\expandafter\def\csname SGAsource@VI.subsec4.1\endcsname{4.1}
\expandafter\def\csname SGAsource@VI.subsec4.12\endcsname{4.12}
\expandafter\def\csname SGAsource@VI.subsec4.9\endcsname{4.9}
\expandafter\def\csname SGAsource@VI.subsec5.9\endcsname{5.9}
\expandafter\def\csname SGAsource@VI.subsec6.0\endcsname{6.0}
\expandafter\def\csname SGAsource@VI.subsec6.10\endcsname{6.10}
\expandafter\def\csname SGAsource@VI.subsec6.11\endcsname{6.11}
\expandafter\def\csname SGAsource@VI.subsec7.1\endcsname{7.1}
\expandafter\def\csname SGAsource@VI.subsec7.2\endcsname{7.2}
\expandafter\def\csname SGAsource@VI.subsec7.3\endcsname{7.3}
\expandafter\def\csname SGAsource@VI.subsec7.4\endcsname{7.4}
\expandafter\def\csname SGAsource@VI.subsec8.1\endcsname{8.1}
\expandafter\def\csname SGAsource@VI.subsec8.2\endcsname{8.2}
\expandafter\def\csname SGAsource@VI.subsec8.3\endcsname{8.3}
\expandafter\def\csname SGAsource@VI.subsec8.4\endcsname{8.4}
\expandafter\def\csname SGAsource@VI.subsec8.5\endcsname{8.5}
\expandafter\def\csname SGAsource@VI.subsec8.6\endcsname{8.6}
\expandafter\def\csname SGAsource@VI.subsec8.7\endcsname{8.7}
\expandafter\def\csname SGAsource@VI.subsubsec1.20.2\endcsname{1.20.2}
\expandafter\def\csname SGAsource@VI.subsubsec1.22.1\endcsname{1.22.1}
\expandafter\def\csname SGAsource@VI.subsubsec1.30.1\endcsname{1.30.1}
\expandafter\def\csname SGAsource@VI.subsubsec1.5.2\endcsname{1.5.2}
\expandafter\def\csname SGAsource@VI.subsubsec1.5.3\endcsname{1.5.3}
\expandafter\def\csname SGAsource@VI.subsubsec1.6.2\endcsname{1.6.2}
\expandafter\def\csname SGAsource@VI.subsubsec1.7.1\endcsname{1.7.1}
\expandafter\def\csname SGAsource@VI.subsubsec1.9.2\endcsname{1.9.2}
\expandafter\def\csname SGAsource@VI.subsubsec1.9.3\endcsname{1.9.3}
\expandafter\def\csname SGAsource@VI.subsubsec1.9.4\endcsname{1.9.4}
\expandafter\def\csname SGAsource@VI.subsubsec2.15.1\endcsname{2.15.1}
\expandafter\def\csname SGAsource@VI.subsubsec2.3.1\endcsname{2.3.1}
\expandafter\def\csname SGAsource@VI.subsubsec2.4.1\endcsname{2.4.1}
\expandafter\def\csname SGAsource@VI.subsubsec2.4.2\endcsname{2.4.2}
\expandafter\def\csname SGAsource@VI.subsubsec2.4.3\endcsname{2.4.3}
\expandafter\def\csname SGAsource@VI.subsubsec2.4.4\endcsname{2.4.4}
\expandafter\def\csname SGAsource@VI.subsubsec2.4.5\endcsname{2.4.5}
\expandafter\def\csname SGAsource@VI.subsubsec2.4.6\endcsname{2.4.6}
\expandafter\def\csname SGAsource@VI.subsubsec2.4.7\endcsname{2.4.7}
\expandafter\def\csname SGAsource@VI.subsubsec2.9.2\endcsname{2.9.2}
\expandafter\def\csname SGAsource@VI.subsubsec3.1.1\endcsname{3.1.1}
\expandafter\def\csname SGAsource@VI.subsubsec3.7.1\endcsname{3.7.1}
\expandafter\def\csname SGAsource@VI.subsubsec4.9.10\endcsname{4.9.10}
\expandafter\def\csname SGAsource@VI.subsubsec4.9.3\endcsname{4.9.3}
\expandafter\def\csname SGAsource@VI.subsubsec6.1.1\endcsname{6.1.1}
\expandafter\def\csname SGAsource@VI.subsubsec6.1.2\endcsname{6.1.2}
\expandafter\def\csname SGAsource@VI.subsubsec6.1.3\endcsname{6.1.3}
\expandafter\def\csname SGAsource@VI.subsubsec6.1.4\endcsname{6.1.4}
\expandafter\def\csname SGAsource@VI.subsubsec6.4.0\endcsname{6.4.0}
\expandafter\def\csname SGAsource@VI.subsubsec7.1.2\endcsname{7.1.2}
\expandafter\def\csname SGAsource@VI.subsubsec7.1.3\endcsname{7.1.3}
\expandafter\def\csname SGAsource@VI.subsubsec7.1.4\endcsname{7.1.4}
\expandafter\def\csname SGAsource@VI.subsubsec7.1.6\endcsname{7.1.6}
\expandafter\def\csname SGAsource@VI.subsubsec7.1.7\endcsname{7.1.7}
\expandafter\def\csname SGAsource@VI.subsubsec7.1.8\endcsname{7.1.8}
\expandafter\def\csname SGAsource@VI.subsubsec7.1.9\endcsname{7.1.9}
\expandafter\def\csname SGAsource@VI.subsubsec7.2.1\endcsname{7.2.1}
\expandafter\def\csname SGAsource@VI.subsubsec7.2.2\endcsname{7.2.2}
\expandafter\def\csname SGAsource@VI.subsubsec7.2.3\endcsname{7.2.3}
\expandafter\def\csname SGAsource@VI.subsubsec7.2.4\endcsname{7.2.4}
\expandafter\def\csname SGAsource@VI.subsubsec7.2.5\endcsname{7.2.5}
\expandafter\def\csname SGAsource@VI.subsubsec7.2.6\endcsname{7.2.6}
\expandafter\def\csname SGAsource@VI.subsubsec7.4.5\endcsname{7.4.5}
\expandafter\def\csname SGAsource@VI.subsubsec7.4.6\endcsname{7.4.6}
\expandafter\def\csname SGAsource@VI.subsubsec7.4.9\endcsname{7.4.9}
\expandafter\def\csname SGAsource@VI.subsubsec8.1.2\endcsname{8.1.2}
\expandafter\def\csname SGAsource@VI.subsubsec8.1.3\endcsname{8.1.3}
\expandafter\def\csname SGAsource@VI.subsubsec8.1.4\endcsname{8.1.4}
\expandafter\def\csname SGAsource@VI.subsubsec8.1.5\endcsname{8.1.5}
\expandafter\def\csname SGAsource@VI.subsubsec8.1.6\endcsname{8.1.6}
\expandafter\def\csname SGAsource@VI.subsubsec8.2.10\endcsname{8.2.10}
\expandafter\def\csname SGAsource@VI.subsubsec8.2.11\endcsname{8.2.11}
\expandafter\def\csname SGAsource@VI.subsubsec8.2.12\endcsname{8.2.12}
\expandafter\def\csname SGAsource@VI.subsubsec8.2.6\endcsname{8.2.6}
\expandafter\def\csname SGAsource@VI.subsubsec8.2.7\endcsname{8.2.7}
\expandafter\def\csname SGAsource@VI.subsubsec8.2.8\endcsname{8.2.8}
\expandafter\def\csname SGAsource@VI.subsubsec8.3.1\endcsname{8.3.1}
\expandafter\def\csname SGAsource@VI.subsubsec8.3.10\endcsname{8.3.10}
\expandafter\def\csname SGAsource@VI.subsubsec8.3.12\endcsname{8.3.12}
\expandafter\def\csname SGAsource@VI.subsubsec8.3.2\endcsname{8.3.2}
\expandafter\def\csname SGAsource@VI.subsubsec8.3.5\endcsname{8.3.5}
\expandafter\def\csname SGAsource@VI.subsubsec8.3.7\endcsname{8.3.7}
\expandafter\def\csname SGAsource@VI.subsubsec8.3.8\endcsname{8.3.8}
\expandafter\def\csname SGAsource@VI.subsubsec8.3.9\endcsname{8.3.9}
\expandafter\def\csname SGAsource@VI.subsubsec8.4.1\endcsname{8.4.1}
\expandafter\def\csname SGAsource@VI.subsubsec8.4.2\endcsname{8.4.2}
\expandafter\def\csname SGAsource@VI.subsubsec8.4.3\endcsname{8.4.3}
\expandafter\def\csname SGAsource@VI.subsubsec8.4.4\endcsname{8.4.4}
\expandafter\def\csname SGAsource@VI.subsubsec8.4.5\endcsname{8.4.5}
\expandafter\def\csname SGAsource@VI.subsubsec8.4.6\endcsname{8.4.6}
\expandafter\def\csname SGAsource@VI.subsubsec8.4.7\endcsname{8.4.7}
\expandafter\def\csname SGAsource@VI.subsubsec8.5.1\endcsname{8.5.1}
\expandafter\def\csname SGAsource@VI.subsubsec8.6.1\endcsname{8.6.1}
\expandafter\def\csname SGAsource@VI.subsubsec8.6.2\endcsname{8.6.2}
\expandafter\def\csname SGAsource@VI.subsubsec8.6.3\endcsname{8.6.3}
\expandafter\def\csname SGAsource@VI.subsubsec8.6.4\endcsname{8.6.4}
\expandafter\def\csname SGAsource@VI.subsubsec8.7.1\endcsname{8.7.1}
\expandafter\def\csname SGAsource@VI.thm1.23\endcsname{1.23}
\expandafter\def\csname SGAsource@VI.thm5.1\endcsname{5.1}
\expandafter\def\csname SGAsource@VI.thm7.4.13.1\endcsname{7.4.13.1}
\expandafter\def\csname SGAsource@VI.thm8.2.3\endcsname{8.2.3}
\expandafter\def\csname SGAsource@VI.thm8.2.9\endcsname{8.2.9}
\expandafter\def\csname SGAsource@VI.thm8.3.13\endcsname{8.3.13}
\expandafter\def\csname SGAsource@VI.thm8.7.3\endcsname{8.7.3}
\expandafter\def\csname SGAsource@VII.cor3.2\endcsname{3.2}
\expandafter\def\csname SGAsource@VII.cor4.4\endcsname{4.4}
\expandafter\def\csname SGAsource@VII.cor5.10\endcsname{5.10}
\expandafter\def\csname SGAsource@VII.cor5.11\endcsname{5.11}
\expandafter\def\csname SGAsource@VII.cor5.12\endcsname{5.12}
\expandafter\def\csname SGAsource@VII.cor5.13\endcsname{5.13}
\expandafter\def\csname SGAsource@VII.cor5.8\endcsname{5.8}
\expandafter\def\csname SGAsource@VII.cor5.9\endcsname{5.9}
\expandafter\def\csname SGAsource@VII.lem5.6\endcsname{5.6}
\expandafter\def\csname SGAsource@VII.prop3.1\endcsname{3.1}
\expandafter\def\csname SGAsource@VII.prop3.3\endcsname{3.3}
\expandafter\def\csname SGAsource@VII.prop4.1\endcsname{4.1}
\expandafter\def\csname SGAsource@VII.prop4.3\endcsname{4.3}
\expandafter\def\csname SGAsource@VII.rem2.1\endcsname{2.1}
\expandafter\def\csname SGAsource@VII.rem4.5\endcsname{4.5}
\expandafter\def\csname SGAsource@VII.rem5.14\endcsname{5.14}
\expandafter\def\csname SGAsource@VII.sec1\endcsname{1}
\expandafter\def\csname SGAsource@VII.sec2\endcsname{2}
\expandafter\def\csname SGAsource@VII.sec3\endcsname{3}
\expandafter\def\csname SGAsource@VII.sec4\endcsname{4}
\expandafter\def\csname SGAsource@VII.sec5\endcsname{5}
\expandafter\def\csname SGAsource@VII.subsec1.1\endcsname{1.1}
\expandafter\def\csname SGAsource@VII.subsec1.2\endcsname{1.2}
\expandafter\def\csname SGAsource@VII.subsec1.3\endcsname{1.3}
\expandafter\def\csname SGAsource@VII.subsec1.4\endcsname{1.4}
\expandafter\def\csname SGAsource@VII.subsec1.5\endcsname{1.5}
\expandafter\def\csname SGAsource@VII.subsec1.6\endcsname{1.6}
\expandafter\def\csname SGAsource@VII.subsec1.7\endcsname{1.7}
\expandafter\def\csname SGAsource@VII.subsec4.0\endcsname{4.0}
\expandafter\def\csname SGAsource@VII.subsec4.2\endcsname{4.2}
\expandafter\def\csname SGAsource@VII.subsec5.1\endcsname{5.1}
\expandafter\def\csname SGAsource@VII.subsec5.2\endcsname{5.2}
\expandafter\def\csname SGAsource@VII.subsec5.3\endcsname{5.3}
\expandafter\def\csname SGAsource@VII.subsec5.4\endcsname{5.4}
\expandafter\def\csname SGAsource@VII.subsec5.5\endcsname{5.5}
\expandafter\def\csname SGAsource@VII.subsubsec4.2.1\endcsname{4.2.1}
\expandafter\def\csname SGAsource@VII.subsubsec4.2.2\endcsname{4.2.2}
\expandafter\def\csname SGAsource@VII.thm5.7\endcsname{5.7}
\expandafter\def\csname SGAsource@VIII.cor1.2\endcsname{1.2}
\expandafter\def\csname SGAsource@VIII.cor2.2\endcsname{2.2}
\expandafter\def\csname SGAsource@VIII.cor2.3\endcsname{2.3}
\expandafter\def\csname SGAsource@VIII.cor2.4\endcsname{2.4}
\expandafter\def\csname SGAsource@VIII.cor2.5\endcsname{2.5}
\expandafter\def\csname SGAsource@VIII.cor3.6\endcsname{3.6}
\expandafter\def\csname SGAsource@VIII.cor3.7\endcsname{3.7}
\expandafter\def\csname SGAsource@VIII.cor3.8\endcsname{3.8}
\expandafter\def\csname SGAsource@VIII.cor4.7\endcsname{4.7}
\expandafter\def\csname SGAsource@VIII.cor4.8\endcsname{4.8}
\expandafter\def\csname SGAsource@VIII.cor5.6\endcsname{5.6}
\expandafter\def\csname SGAsource@VIII.cor5.7\endcsname{5.7}
\expandafter\def\csname SGAsource@VIII.cor6.2\endcsname{6.2}
\expandafter\def\csname SGAsource@VIII.cor6.4\endcsname{6.4}
\expandafter\def\csname SGAsource@VIII.cor7.5\endcsname{7.5}
\expandafter\def\csname SGAsource@VIII.cor7.6\endcsname{7.6}
\expandafter\def\csname SGAsource@VIII.cor8.5\endcsname{8.5}
\expandafter\def\csname SGAsource@VIII.cor8.6\endcsname{8.6}
\expandafter\def\csname SGAsource@VIII.cor9.2\endcsname{9.2}
\expandafter\def\csname SGAsource@VIII.def3.3\endcsname{3.3}
\expandafter\def\csname SGAsource@VIII.def4.2\endcsname{4.2}
\expandafter\def\csname SGAsource@VIII.def4.3\endcsname{4.3}
\expandafter\def\csname SGAsource@VIII.def7.2\endcsname{7.2}
\expandafter\def\csname SGAsource@VIII.eq5.1.1\endcsname{5.1.1}
\expandafter\def\csname SGAsource@VIII.eq5.1.2\endcsname{5.1.2}
\expandafter\def\csname SGAsource@VIII.eq5.4.1\endcsname{5.4.1}
\expandafter\def\csname SGAsource@VIII.eq5.4.2\endcsname{5.4.2}
\expandafter\def\csname SGAsource@VIII.exem1.3\endcsname{1.3}
\expandafter\def\csname SGAsource@VIII.lem7.3\endcsname{7.3}
\expandafter\def\csname SGAsource@VIII.lem8.2\endcsname{8.2}
\expandafter\def\csname SGAsource@VIII.lem9.3\endcsname{9.3}
\expandafter\def\csname SGAsource@VIII.prop2.1\endcsname{2.1}
\expandafter\def\csname SGAsource@VIII.prop3.9\endcsname{3.9}
\expandafter\def\csname SGAsource@VIII.prop4.6\endcsname{4.6}
\expandafter\def\csname SGAsource@VIII.prop5.5\endcsname{5.5}
\expandafter\def\csname SGAsource@VIII.prop6.1\endcsname{6.1}
\expandafter\def\csname SGAsource@VIII.prop7.4\endcsname{7.4}
\expandafter\def\csname SGAsource@VIII.prop8.1\endcsname{8.1}
\expandafter\def\csname SGAsource@VIII.prop8.4\endcsname{8.4}
\expandafter\def\csname SGAsource@VIII.prop9.1\endcsname{9.1}
\expandafter\def\csname SGAsource@VIII.rem1.4\endcsname{1.4}
\expandafter\def\csname SGAsource@VIII.rem3.13\endcsname{3.13}
\expandafter\def\csname SGAsource@VIII.rem3.2\endcsname{3.2}
\expandafter\def\csname SGAsource@VIII.rem5.3\endcsname{5.3}
\expandafter\def\csname SGAsource@VIII.rem5.8\endcsname{5.8}
\expandafter\def\csname SGAsource@VIII.rem8.3\endcsname{8.3}
\expandafter\def\csname SGAsource@VIII.sch3.10\endcsname{3.10}
\expandafter\def\csname SGAsource@VIII.sec1\endcsname{1}
\expandafter\def\csname SGAsource@VIII.sec2\endcsname{2}
\expandafter\def\csname SGAsource@VIII.sec3\endcsname{3}
\expandafter\def\csname SGAsource@VIII.sec4\endcsname{4}
\expandafter\def\csname SGAsource@VIII.sec5\endcsname{5}
\expandafter\def\csname SGAsource@VIII.sec6\endcsname{6}
\expandafter\def\csname SGAsource@VIII.sec7\endcsname{7}
\expandafter\def\csname SGAsource@VIII.sec8\endcsname{8}
\expandafter\def\csname SGAsource@VIII.sec9\endcsname{9}
\expandafter\def\csname SGAsource@VIII.subsec3.1\endcsname{3.1}
\expandafter\def\csname SGAsource@VIII.subsec3.11\endcsname{3.11}
\expandafter\def\csname SGAsource@VIII.subsec3.12\endcsname{3.12}
\expandafter\def\csname SGAsource@VIII.subsec3.4\endcsname{3.4}
\expandafter\def\csname SGAsource@VIII.subsec4.1\endcsname{4.1}
\expandafter\def\csname SGAsource@VIII.subsec4.4\endcsname{4.4}
\expandafter\def\csname SGAsource@VIII.subsec4.5\endcsname{4.5}
\expandafter\def\csname SGAsource@VIII.subsec5.1\endcsname{5.1}
\expandafter\def\csname SGAsource@VIII.subsec5.4\endcsname{5.4}
\expandafter\def\csname SGAsource@VIII.subsec6.5\endcsname{6.5}
\expandafter\def\csname SGAsource@VIII.subsec6.6\endcsname{6.6}
\expandafter\def\csname SGAsource@VIII.subsec7.1\endcsname{7.1}
\expandafter\def\csname SGAsource@VIII.subsec7.7\endcsname{7.7}
\expandafter\def\csname SGAsource@VIII.subsec7.8\endcsname{7.8}
\expandafter\def\csname SGAsource@VIII.subsubsec7.9.1\endcsname{7.9.1}
\expandafter\def\csname SGAsource@VIII.subsubsec7.9.2\endcsname{7.9.2}
\expandafter\def\csname SGAsource@VIII.subsubsec7.9.3\endcsname{7.9.3}
\expandafter\def\csname SGAsource@VIII.subsubsec9.4.1\endcsname{9.4.1}
\expandafter\def\csname SGAsource@VIII.subsubsec9.4.2\endcsname{9.4.2}
\expandafter\def\csname SGAsource@VIII.subsubsec9.4.3\endcsname{9.4.3}
\expandafter\def\csname SGAsource@VIII.subsubsec9.4.4\endcsname{9.4.4}
\expandafter\def\csname SGAsource@VIII.thm1.1\endcsname{1.1}
\expandafter\def\csname SGAsource@VIII.thm3.5\endcsname{3.5}
\expandafter\def\csname SGAsource@VIII.thm5.2\endcsname{5.2}
\expandafter\def\csname SGAsource@VIII.thm6.3\endcsname{6.3}
\expandafter\def\csname SGAsource@VIII.thm7.9\endcsname{7.9}
\expandafter\def\csname SGAsource@VIII.thm9.4\endcsname{9.4}
\expandafter\def\csname SGAsource@Vbis.cor1.2.11\endcsname{1.2.11}
\expandafter\def\csname SGAsource@Vbis.cor1.3.11\endcsname{1.3.11}
\expandafter\def\csname SGAsource@Vbis.cor1.3.12\endcsname{1.3.12}
\expandafter\def\csname SGAsource@Vbis.cor2.3.6\endcsname{2.3.6}
\expandafter\def\csname SGAsource@Vbis.cor2.3.7\endcsname{2.3.7}
\expandafter\def\csname SGAsource@Vbis.cor2.4.14\endcsname{2.4.14}
\expandafter\def\csname SGAsource@Vbis.cor3.1.7\endcsname{3.1.7}
\expandafter\def\csname SGAsource@Vbis.cor3.3.5\endcsname{3.3.5}
\expandafter\def\csname SGAsource@Vbis.cor4.1.2\endcsname{4.1.2}
\expandafter\def\csname SGAsource@Vbis.cor4.1.6\endcsname{4.1.6}
\expandafter\def\csname SGAsource@Vbis.cor4.1.7\endcsname{4.1.7}
\expandafter\def\csname SGAsource@Vbis.cor4.1.9\endcsname{4.1.9}
\expandafter\def\csname SGAsource@Vbis.cor4.3.5\endcsname{4.3.5}
\expandafter\def\csname SGAsource@Vbis.def1.2.1\endcsname{1.2.1}
\expandafter\def\csname SGAsource@Vbis.def1.2.14\endcsname{1.2.14}
\expandafter\def\csname SGAsource@Vbis.def1.2.2\endcsname{1.2.2}
\expandafter\def\csname SGAsource@Vbis.def1.3.1\endcsname{1.3.1}
\expandafter\def\csname SGAsource@Vbis.def1.3.2\endcsname{1.3.2}
\expandafter\def\csname SGAsource@Vbis.def2.1.2\endcsname{2.1.2}
\expandafter\def\csname SGAsource@Vbis.def2.2.2\endcsname{2.2.2}
\expandafter\def\csname SGAsource@Vbis.def2.2.4\endcsname{2.2.4}
\expandafter\def\csname SGAsource@Vbis.def2.2.6\endcsname{2.2.6}
\expandafter\def\csname SGAsource@Vbis.def2.4.0.1\endcsname{2.4.0.1}
\expandafter\def\csname SGAsource@Vbis.def2.4.10\endcsname{2.4.10}
\expandafter\def\csname SGAsource@Vbis.def2.4.11\endcsname{2.4.11}
\expandafter\def\csname SGAsource@Vbis.def2.4.15\endcsname{2.4.15}
\expandafter\def\csname SGAsource@Vbis.def2.4.16\endcsname{2.4.16}
\expandafter\def\csname SGAsource@Vbis.def2.4.3.1\endcsname{2.4.3.1}
\expandafter\def\csname SGAsource@Vbis.def2.4.8\endcsname{2.4.8}
\expandafter\def\csname SGAsource@Vbis.def3.1.2\endcsname{3.1.2}
\expandafter\def\csname SGAsource@Vbis.def3.2.2.4\endcsname{3.2.2.4}
\expandafter\def\csname SGAsource@Vbis.def5.1.1\endcsname{5.1.1}
\expandafter\def\csname SGAsource@Vbis.def5.1.5\endcsname{5.1.5}
\expandafter\def\csname SGAsource@Vbis.def5.2.6\endcsname{5.2.6}
\expandafter\def\csname SGAsource@Vbis.eq1.1\endcsname{1.1}
\expandafter\def\csname SGAsource@Vbis.eq1.2\endcsname{1.2}
\expandafter\def\csname SGAsource@Vbis.eq1.3\endcsname{1.3}
\expandafter\def\csname SGAsource@Vbis.eq1.3.10.2\endcsname{1.3.10.2}
\expandafter\def\csname SGAsource@Vbis.eq1.3.10.3\endcsname{1.3.10.3}
\expandafter\def\csname SGAsource@Vbis.eq1.3.10.4\endcsname{1.3.10.4}
\expandafter\def\csname SGAsource@Vbis.eq2.1\endcsname{2.1}
\expandafter\def\csname SGAsource@Vbis.eq2.2\endcsname{2.2}
\expandafter\def\csname SGAsource@Vbis.eq2.3\endcsname{2.3}
\expandafter\def\csname SGAsource@Vbis.eq2.3.10.2\endcsname{2.3.10.2}
\expandafter\def\csname SGAsource@Vbis.eq2.3.2.1\endcsname{2.3.2.1}
\expandafter\def\csname SGAsource@Vbis.eq2.3.3.1\endcsname{2.3.3.1}
\expandafter\def\csname SGAsource@Vbis.eq2.3.3.2\endcsname{2.3.3.2}
\expandafter\def\csname SGAsource@Vbis.eq2.3.4.1\endcsname{2.3.4.1}
\expandafter\def\csname SGAsource@Vbis.eq2.3.4.2\endcsname{2.3.4.2}
\expandafter\def\csname SGAsource@Vbis.eq2.3.4.3\endcsname{2.3.4.3}
\expandafter\def\csname SGAsource@Vbis.eq2.3.5.1\endcsname{2.3.5.1}
\expandafter\def\csname SGAsource@Vbis.eq2.3.5.2\endcsname{2.3.5.2}
\expandafter\def\csname SGAsource@Vbis.eq2.3.9.1\endcsname{2.3.9.1}
\expandafter\def\csname SGAsource@Vbis.eq2.4\endcsname{2.4}
\expandafter\def\csname SGAsource@Vbis.eq2.4.4.1\endcsname{2.4.4.1}
\expandafter\def\csname SGAsource@Vbis.eq2.5.1.1\endcsname{2.5.1.1}
\expandafter\def\csname SGAsource@Vbis.eq2.5.1.2\endcsname{2.5.1.2}
\expandafter\def\csname SGAsource@Vbis.eq2.5.2.1\endcsname{2.5.2.1}
\expandafter\def\csname SGAsource@Vbis.eq2.5.3.1\endcsname{2.5.3.1}
\expandafter\def\csname SGAsource@Vbis.eq2.5.3.2\endcsname{2.5.3.2}
\expandafter\def\csname SGAsource@Vbis.eq2.5.3.3\endcsname{2.5.3.3}
\expandafter\def\csname SGAsource@Vbis.eq2.5.5.1\endcsname{2.5.5.1}
\expandafter\def\csname SGAsource@Vbis.eq2.5.5.2\endcsname{2.5.5.2}
\expandafter\def\csname SGAsource@Vbis.eq2.5.5.3\endcsname{2.5.5.3}
\expandafter\def\csname SGAsource@Vbis.eq3.0.1.1\endcsname{3.0.1.1}
\expandafter\def\csname SGAsource@Vbis.eq3.0.1.2\endcsname{3.0.1.2}
\expandafter\def\csname SGAsource@Vbis.eq3.0.2.1\endcsname{3.0.2.1}
\expandafter\def\csname SGAsource@Vbis.eq3.1.1.1\endcsname{3.1.1.1}
\expandafter\def\csname SGAsource@Vbis.eq3.1.1.2\endcsname{3.1.1.2}
\expandafter\def\csname SGAsource@Vbis.eq3.1.1.3\endcsname{3.1.1.3}
\expandafter\def\csname SGAsource@Vbis.eq3.1.1.4\endcsname{3.1.1.4}
\expandafter\def\csname SGAsource@Vbis.eq3.1.1.5\endcsname{3.1.1.5}
\expandafter\def\csname SGAsource@Vbis.eq3.1.1.6\endcsname{3.1.1.6}
\expandafter\def\csname SGAsource@Vbis.eq3.1.1.7\endcsname{3.1.1.7}
\expandafter\def\csname SGAsource@Vbis.eq3.1.1.8\endcsname{3.1.1.8}
\expandafter\def\csname SGAsource@Vbis.eq3.2\endcsname{3.2}
\expandafter\def\csname SGAsource@Vbis.eq3.2.2.1\endcsname{3.2.2.1}
\expandafter\def\csname SGAsource@Vbis.eq3.2.2.2\endcsname{3.2.2.2}
\expandafter\def\csname SGAsource@Vbis.eq3.2.2.3\endcsname{3.2.2.3}
\expandafter\def\csname SGAsource@Vbis.eq3.2.3.1\endcsname{3.2.3.1}
\expandafter\def\csname SGAsource@Vbis.eq3.2.3.2\endcsname{3.2.3.2}
\expandafter\def\csname SGAsource@Vbis.eq3.3\endcsname{3.3}
\expandafter\def\csname SGAsource@Vbis.eq4.1\endcsname{4.1}
\expandafter\def\csname SGAsource@Vbis.eq4.1.3.1\endcsname{4.1.3.1}
\expandafter\def\csname SGAsource@Vbis.lem1.2.10\endcsname{1.2.10}
\expandafter\def\csname SGAsource@Vbis.lem1.2.10.1\endcsname{1.2.10.1}
\expandafter\def\csname SGAsource@Vbis.lem1.2.16\endcsname{1.2.16}
\expandafter\def\csname SGAsource@Vbis.lem1.3.10.1\endcsname{1.3.10.1}
\expandafter\def\csname SGAsource@Vbis.lem1.3.10.5\endcsname{1.3.10.5}
\expandafter\def\csname SGAsource@Vbis.lem2.2.3.1\endcsname{2.2.3.1}
\expandafter\def\csname SGAsource@Vbis.lem2.3.2.2\endcsname{2.3.2.2}
\expandafter\def\csname SGAsource@Vbis.lem2.4.3.2\endcsname{2.4.3.2}
\expandafter\def\csname SGAsource@Vbis.lem2.4.4.2\endcsname{2.4.4.2}
\expandafter\def\csname SGAsource@Vbis.lem2.4.5.1\endcsname{2.4.5.1}
\expandafter\def\csname SGAsource@Vbis.lem2.4.7.1\endcsname{2.4.7.1}
\expandafter\def\csname SGAsource@Vbis.lem3.0.2.3\endcsname{3.0.2.3}
\expandafter\def\csname SGAsource@Vbis.lem3.0.2.4\endcsname{3.0.2.4}
\expandafter\def\csname SGAsource@Vbis.lem3.1.1.9\endcsname{3.1.1.9}
\expandafter\def\csname SGAsource@Vbis.lem3.1.5.1\endcsname{3.1.5.1}
\expandafter\def\csname SGAsource@Vbis.lem3.2.3.3\endcsname{3.2.3.3}
\expandafter\def\csname SGAsource@Vbis.lem3.3.3.1\endcsname{3.3.3.1}
\expandafter\def\csname SGAsource@Vbis.lem3.3.3.2\endcsname{3.3.3.2}
\expandafter\def\csname SGAsource@Vbis.lem3.3.3.3\endcsname{3.3.3.3}
\expandafter\def\csname SGAsource@Vbis.lem3.3.4.1\endcsname{3.3.4.1}
\expandafter\def\csname SGAsource@Vbis.lem3.3.4.2\endcsname{3.3.4.2}
\expandafter\def\csname SGAsource@Vbis.lem4.1.3\endcsname{4.1.3}
\expandafter\def\csname SGAsource@Vbis.lem4.1.4\endcsname{4.1.4}
\expandafter\def\csname SGAsource@Vbis.lem4.1.5\endcsname{4.1.5}
\expandafter\def\csname SGAsource@Vbis.lem5.3.2\endcsname{5.3.2}
\expandafter\def\csname SGAsource@Vbis.par1.3.6.1\endcsname{1.3.6.1}
\expandafter\def\csname SGAsource@Vbis.par2.3.1.1\endcsname{2.3.1.1}
\expandafter\def\csname SGAsource@Vbis.par2.3.1.2\endcsname{2.3.1.2}
\expandafter\def\csname SGAsource@Vbis.par2.3.1.3\endcsname{2.3.1.3}
\expandafter\def\csname SGAsource@Vbis.par2.4.1.1\endcsname{2.4.1.1}
\expandafter\def\csname SGAsource@Vbis.par3.0.2.2\endcsname{3.0.2.2}
\expandafter\def\csname SGAsource@Vbis.prop1.2.12\endcsname{1.2.12}
\expandafter\def\csname SGAsource@Vbis.prop1.2.15\endcsname{1.2.15}
\expandafter\def\csname SGAsource@Vbis.prop1.2.9\endcsname{1.2.9}
\expandafter\def\csname SGAsource@Vbis.prop1.3.10\endcsname{1.3.10}
\expandafter\def\csname SGAsource@Vbis.prop1.3.7\endcsname{1.3.7}
\expandafter\def\csname SGAsource@Vbis.prop2.2.7\endcsname{2.2.7}
\expandafter\def\csname SGAsource@Vbis.prop2.3.5\endcsname{2.3.5}
\expandafter\def\csname SGAsource@Vbis.prop2.4.13\endcsname{2.4.13}
\expandafter\def\csname SGAsource@Vbis.prop2.4.6\endcsname{2.4.6}
\expandafter\def\csname SGAsource@Vbis.prop2.5.4\endcsname{2.5.4}
\expandafter\def\csname SGAsource@Vbis.prop2.5.5\endcsname{2.5.5}
\expandafter\def\csname SGAsource@Vbis.prop2.5.6\endcsname{2.5.6}
\expandafter\def\csname SGAsource@Vbis.prop2.5.7\endcsname{2.5.7}
\expandafter\def\csname SGAsource@Vbis.prop3.1.1\endcsname{3.1.1}
\expandafter\def\csname SGAsource@Vbis.prop3.1.4\endcsname{3.1.4}
\expandafter\def\csname SGAsource@Vbis.prop3.1.5\endcsname{3.1.5}
\expandafter\def\csname SGAsource@Vbis.prop3.1.6\endcsname{3.1.6}
\expandafter\def\csname SGAsource@Vbis.prop3.2.1\endcsname{3.2.1}
\expandafter\def\csname SGAsource@Vbis.prop3.2.4\endcsname{3.2.4}
\expandafter\def\csname SGAsource@Vbis.prop3.3.1\endcsname{3.3.1}
\expandafter\def\csname SGAsource@Vbis.prop3.3.4\endcsname{3.3.4}
\expandafter\def\csname SGAsource@Vbis.prop4.1.8\endcsname{4.1.8}
\expandafter\def\csname SGAsource@Vbis.prop4.2.1\endcsname{4.2.1}
\expandafter\def\csname SGAsource@Vbis.prop4.3.2\endcsname{4.3.2}
\expandafter\def\csname SGAsource@Vbis.prop4.3.3\endcsname{4.3.3}
\expandafter\def\csname SGAsource@Vbis.prop5.1.3\endcsname{5.1.3}
\expandafter\def\csname SGAsource@Vbis.prop5.1.7\endcsname{5.1.7}
\expandafter\def\csname SGAsource@Vbis.prop5.1.8\endcsname{5.1.8}
\expandafter\def\csname SGAsource@Vbis.prop5.2.3\endcsname{5.2.3}
\expandafter\def\csname SGAsource@Vbis.prop5.2.5\endcsname{5.2.5}
\expandafter\def\csname SGAsource@Vbis.prop5.2.8\endcsname{5.2.8}
\expandafter\def\csname SGAsource@Vbis.prop5.3.4\endcsname{5.3.4}
\expandafter\def\csname SGAsource@Vbis.rem1.3.13\endcsname{1.3.13}
\expandafter\def\csname SGAsource@Vbis.rem1.3.3\endcsname{1.3.3}
\expandafter\def\csname SGAsource@Vbis.rem2.2.5\endcsname{2.2.5}
\expandafter\def\csname SGAsource@Vbis.rem2.3.8\endcsname{2.3.8}
\expandafter\def\csname SGAsource@Vbis.rem2.4.9\endcsname{2.4.9}
\expandafter\def\csname SGAsource@Vbis.rem4.1.10\endcsname{4.1.10}
\expandafter\def\csname SGAsource@Vbis.rem4.3.4\endcsname{4.3.4}
\expandafter\def\csname SGAsource@Vbis.rem5.2.9\endcsname{5.2.9}
\expandafter\def\csname SGAsource@Vbis.sch1.2.17\endcsname{1.2.17}
\expandafter\def\csname SGAsource@Vbis.sec1\endcsname{1}
\expandafter\def\csname SGAsource@Vbis.sec2\endcsname{2}
\expandafter\def\csname SGAsource@Vbis.sec3\endcsname{3}
\expandafter\def\csname SGAsource@Vbis.sec4\endcsname{4}
\expandafter\def\csname SGAsource@Vbis.sec5\endcsname{5}
\expandafter\def\csname SGAsource@Vbis.subsec1.1\endcsname{1.1}
\expandafter\def\csname SGAsource@Vbis.subsec1.2\endcsname{1.2}
\expandafter\def\csname SGAsource@Vbis.subsec1.3\endcsname{1.3}
\expandafter\def\csname SGAsource@Vbis.subsec2.1\endcsname{2.1}
\expandafter\def\csname SGAsource@Vbis.subsec2.2\endcsname{2.2}
\expandafter\def\csname SGAsource@Vbis.subsec2.3\endcsname{2.3}
\expandafter\def\csname SGAsource@Vbis.subsec2.4\endcsname{2.4}
\expandafter\def\csname SGAsource@Vbis.subsec2.5\endcsname{2.5}
\expandafter\def\csname SGAsource@Vbis.subsec3.0\endcsname{3.0}
\expandafter\def\csname SGAsource@Vbis.subsec3.1\endcsname{3.1}
\expandafter\def\csname SGAsource@Vbis.subsec3.2\endcsname{3.2}
\expandafter\def\csname SGAsource@Vbis.subsec3.3\endcsname{3.3}
\expandafter\def\csname SGAsource@Vbis.subsec4.1\endcsname{4.1}
\expandafter\def\csname SGAsource@Vbis.subsec4.2\endcsname{4.2}
\expandafter\def\csname SGAsource@Vbis.subsec4.3\endcsname{4.3}
\expandafter\def\csname SGAsource@Vbis.subsec5.1\endcsname{5.1}
\expandafter\def\csname SGAsource@Vbis.subsec5.2\endcsname{5.2}
\expandafter\def\csname SGAsource@Vbis.subsec5.3\endcsname{5.3}
\expandafter\def\csname SGAsource@Vbis.subsubsec1.1.1\endcsname{1.1.1}
\expandafter\def\csname SGAsource@Vbis.subsubsec1.1.2\endcsname{1.1.2}
\expandafter\def\csname SGAsource@Vbis.subsubsec1.1.3\endcsname{1.1.3}
\expandafter\def\csname SGAsource@Vbis.subsubsec1.2.13\endcsname{1.2.13}
\expandafter\def\csname SGAsource@Vbis.subsubsec1.2.4\endcsname{1.2.4}
\expandafter\def\csname SGAsource@Vbis.subsubsec1.2.5\endcsname{1.2.5}
\expandafter\def\csname SGAsource@Vbis.subsubsec1.2.6\endcsname{1.2.6}
\expandafter\def\csname SGAsource@Vbis.subsubsec1.2.7\endcsname{1.2.7}
\expandafter\def\csname SGAsource@Vbis.subsubsec1.2.8\endcsname{1.2.8}
\expandafter\def\csname SGAsource@Vbis.subsubsec1.3.4\endcsname{1.3.4}
\expandafter\def\csname SGAsource@Vbis.subsubsec1.3.5\endcsname{1.3.5}
\expandafter\def\csname SGAsource@Vbis.subsubsec1.3.6\endcsname{1.3.6}
\expandafter\def\csname SGAsource@Vbis.subsubsec1.3.8\endcsname{1.3.8}
\expandafter\def\csname SGAsource@Vbis.subsubsec1.3.9\endcsname{1.3.9}
\expandafter\def\csname SGAsource@Vbis.subsubsec2.1.1\endcsname{2.1.1}
\expandafter\def\csname SGAsource@Vbis.subsubsec2.1.3\endcsname{2.1.3}
\expandafter\def\csname SGAsource@Vbis.subsubsec2.2.1\endcsname{2.2.1}
\expandafter\def\csname SGAsource@Vbis.subsubsec2.2.3\endcsname{2.2.3}
\expandafter\def\csname SGAsource@Vbis.subsubsec2.3.1\endcsname{2.3.1}
\expandafter\def\csname SGAsource@Vbis.subsubsec2.3.10\endcsname{2.3.10}
\expandafter\def\csname SGAsource@Vbis.subsubsec2.3.11\endcsname{2.3.11}
\expandafter\def\csname SGAsource@Vbis.subsubsec2.3.2\endcsname{2.3.2}
\expandafter\def\csname SGAsource@Vbis.subsubsec2.3.3\endcsname{2.3.3}
\expandafter\def\csname SGAsource@Vbis.subsubsec2.3.4\endcsname{2.3.4}
\expandafter\def\csname SGAsource@Vbis.subsubsec2.3.9\endcsname{2.3.9}
\expandafter\def\csname SGAsource@Vbis.subsubsec2.4.0\endcsname{2.4.0}
\expandafter\def\csname SGAsource@Vbis.subsubsec2.4.1\endcsname{2.4.1}
\expandafter\def\csname SGAsource@Vbis.subsubsec2.4.3\endcsname{2.4.3}
\expandafter\def\csname SGAsource@Vbis.subsubsec2.4.4\endcsname{2.4.4}
\expandafter\def\csname SGAsource@Vbis.subsubsec2.4.5\endcsname{2.4.5}
\expandafter\def\csname SGAsource@Vbis.subsubsec2.4.7\endcsname{2.4.7}
\expandafter\def\csname SGAsource@Vbis.subsubsec2.5.1\endcsname{2.5.1}
\expandafter\def\csname SGAsource@Vbis.subsubsec2.5.2\endcsname{2.5.2}
\expandafter\def\csname SGAsource@Vbis.subsubsec3.0.0\endcsname{3.0.0}
\expandafter\def\csname SGAsource@Vbis.subsubsec3.0.1\endcsname{3.0.1}
\expandafter\def\csname SGAsource@Vbis.subsubsec3.0.2\endcsname{3.0.2}
\expandafter\def\csname SGAsource@Vbis.subsubsec3.1.3\endcsname{3.1.3}
\expandafter\def\csname SGAsource@Vbis.subsubsec3.2.2\endcsname{3.2.2}
\expandafter\def\csname SGAsource@Vbis.subsubsec3.2.3\endcsname{3.2.3}
\expandafter\def\csname SGAsource@Vbis.subsubsec3.3.2\endcsname{3.3.2}
\expandafter\def\csname SGAsource@Vbis.subsubsec4.1.0\endcsname{4.1.0}
\expandafter\def\csname SGAsource@Vbis.subsubsec4.2.0\endcsname{4.2.0}
\expandafter\def\csname SGAsource@Vbis.subsubsec4.3.0\endcsname{4.3.0}
\expandafter\def\csname SGAsource@Vbis.subsubsec4.3.1\endcsname{4.3.1}
\expandafter\def\csname SGAsource@Vbis.subsubsec5.1.0\endcsname{5.1.0}
\expandafter\def\csname SGAsource@Vbis.subsubsec5.1.2\endcsname{5.1.2}
\expandafter\def\csname SGAsource@Vbis.subsubsec5.1.4\endcsname{5.1.4}
\expandafter\def\csname SGAsource@Vbis.subsubsec5.1.6\endcsname{5.1.6}
\expandafter\def\csname SGAsource@Vbis.subsubsec5.2.0\endcsname{5.2.0}
\expandafter\def\csname SGAsource@Vbis.subsubsec5.2.1\endcsname{5.2.1}
\expandafter\def\csname SGAsource@Vbis.subsubsec5.2.2\endcsname{5.2.2}
\expandafter\def\csname SGAsource@Vbis.subsubsec5.2.4\endcsname{5.2.4}
\expandafter\def\csname SGAsource@Vbis.subsubsec5.2.7\endcsname{5.2.7}
\expandafter\def\csname SGAsource@Vbis.subsubsec5.3.1\endcsname{5.3.1}
\expandafter\def\csname SGAsource@Vbis.subsubsec5.3.3\endcsname{5.3.3}
\expandafter\def\csname SGAsource@Vbis.susubsec2.4.2\endcsname{2.4.2}
\expandafter\def\csname SGAsource@Vbis.thm2.4.12\endcsname{2.4.12}
\expandafter\def\csname SGAsource@Vbis.thm3.3.3\endcsname{3.3.3}
\expandafter\def\csname SGAsource@Vbis.thm4.1.1\endcsname{4.1.1}
\expandafter\def\csname SGAsource@page112\endcsname{112}
\expandafter\def\csname SGAsource@page133\endcsname{133}
\expandafter\def\csname SGAsource@page6\endcsname{6}
\expandafter\def\csname SGAsource@page7\endcsname{7}
\expandafter\def\csname SGAsource@sga4.exp.I\endcsname{I}
\expandafter\def\csname SGAsource@sga4.exp.II\endcsname{II}
\expandafter\def\csname SGAsource@sga4.exp.III\endcsname{III}
\expandafter\def\csname SGAsource@sga4.exp.III.terminology.5.3\endcsname{3}
\expandafter\def\csname SGAsource@sga4.exp.IV\endcsname{IV}
\expandafter\def\csname SGAsource@sga4.exp.IX\endcsname{IX}
\expandafter\def\csname SGAsource@sga4.exp.V\endcsname{V}
\expandafter\def\csname SGAsource@sga4.exp.VI\endcsname{VI}
\expandafter\def\csname SGAsource@sga4.exp.VII\endcsname{VII}
\expandafter\def\csname SGAsource@sga4.exp.VIII\endcsname{VIII}
\expandafter\def\csname SGAsource@sga4.exp.Vbis\endcsname{Vbis}
\expandafter\def\csname SGAsource@sga4[1]\endcsname{sga4[1]}
\let\SGAoriginalhyperref\hyperref
\renewcommand{\hyperref}[2][]{\ifcsname SGAsource@#1\endcsname #2\else\SGAoriginalhyperref[#1]{#2}\fi}
\AtBeginDocument{\let\SGAoriginalref\ref\renewcommand{\ref}[1]{\ifcsname SGAsource@#1\endcsname\csname SGAsource@#1\endcsname\else\SGAoriginalref{#1}\fi}}


% Keep editorial note numbers local. The inherited macros let thefootnote
% escape, making later original notes [1], [2], [3] all show the last cnote.
\makeatletter
\renewcommand{\nde}[1]{\ifthenelse{\boolean{original}}{}{%
  \stepcounter{cnote}\begingroup
  \def\thefootnote{\arabic{cnote}}%
  \protect\footnote[\value{cnote}]{编者注: #1}\endgroup}}
\renewcommand{\ndemark}{\ifthenelse{\boolean{original}}{}{%
  \stepcounter{cnote}\begingroup
  \def\thefootnote{\arabic{cnote}}%
  \protect\footnotemark[\value{cnote}]\endgroup}}
\renewcommand{\ndetext}[1]{\ifthenelse{\boolean{original}}{}{%
  \begingroup\def\thefootnote{\arabic{cnote}}%
  \protect\footnotetext[\value{cnote}]{编者注: #1}\endgroup}}
% The two inherited repeated-author-note calls both refer to original note i.
% Explicit note numbers do not advance the ordinary footnote counter.
\renewcommand{\sgarepeatedfootnotemark}[2][1]{%
  \begingroup\c@footnote=#1\relax
  \edef\@thefnmark{\roman{footnote}}%
  \hyperlink{#2}{\@makefnmark}\endgroup}
\makeatother
\renewcommand{\nderef}[1]{第~\ref{#1}~条编者注（第~\pageref{#1}~页）}
